

\documentclass[11pt]{article}

\usepackage[margin=0.85in]{geometry}
\usepackage{amsmath,amssymb,mathtools}
\usepackage{enumitem}
\usepackage{xcolor}
\usepackage[most]{tcolorbox}
\usepackage{fancyhdr}
\usepackage{lastpage}
\usepackage{hyperref}
\usepackage{microtype}
\usepackage{lmodern}

%%%%%%%%%%%%%%%%%%%%%%%%%%%%%%%%%%%%%
% Colors
%%%%%%%%%%%%%%%%%%%%%%%%%%%%%%%%%%%%%

\definecolor{uwpurple}{HTML}{4B2E83}
\definecolor{uwgold}{HTML}{B7A57A}
\definecolor{softpurple}{HTML}{F4F1FA}

%%%%%%%%%%%%%%%%%%%%%%%%%%%%%%%%%%%%%
% Hyperlinks
%%%%%%%%%%%%%%%%%%%%%%%%%%%%%%%%%%%%%

\hypersetup{
    colorlinks=true,
    linkcolor=uwpurple,
    urlcolor=uwpurple,
    citecolor=uwpurple
}

%%%%%%%%%%%%%%%%%%%%%%%%%%%%%%%%%%%%%
% Page Style
%%%%%%%%%%%%%%%%%%%%%%%%%%%%%%%%%%%%%

\pagestyle{fancy}
\fancyhf{}
\lhead{\textbf{MATH 394: Probability I}}
\rhead{\textbf{Practice Problems}}
\cfoot{\thepage\ of \pageref{LastPage}}

\setlength{\headheight}{14pt}
\setlength{\parindent}{0pt}
\setlength{\parskip}{0.45em}

\setlist[enumerate]{
    itemsep=0.35em,
    topsep=0.35em
}

%%%%%%%%%%%%%%%%%%%%%%%%%%%%%%%%%%%%%
% Course Information
%%%%%%%%%%%%%%%%%%%%%%%%%%%%%%%%%%%%%

\newcommand{\course}{MATH 394: Probability I}
\newcommand{\documenttitle}{Question Bank with Solutions \\ (Final)}
\newcommand{\instructor}{Arman Jahangiri}
\newcommand{\department}{Department of Mathematics}
\newcommand{\university}{University of Washington}
\newcommand{\term}{Summer 2026}

%%%%%%%%%%%%%%%%%%%%%%%%%%%%%%%%%%%%%
% Problem Formatting
%%%%%%%%%%%%%%%%%%%%%%%%%%%%%%%%%%%%%

\newcounter{problem}

\newtcolorbox{problemheader}{
    colback=softpurple,
    colframe=uwpurple,
    boxrule=0.8pt,
    arc=2mm,
    left=2mm,
    right=2mm,
    top=1.5mm,
    bottom=1.5mm,
    before skip=1em,
    after skip=0.5em
}

\newenvironment{problem}[1][]{
    \refstepcounter{problem}
    \begin{problemheader}
        \textbf{Problem \theproblem%
        \if\relax\detokenize{#1}\relax
        \else. #1%
        \fi}
    \end{problemheader}
}{
    \vspace{0.5em}
}

%%%%%%%%%%%%%%%%%%%%%%%%%%%%%%%%%%%%%
% Solution Formatting
%%%%%%%%%%%%%%%%%%%%%%%%%%%%%%%%%%%%%

\newenvironment{solution}{
    \par
    \begin{tcolorbox}[
        colback=white,
        colframe=uwgold,
        boxrule=0.7pt,
        arc=2mm,
        left=3mm,
        right=3mm,
        top=2mm,
        bottom=2mm,
        breakable,
        before skip=0.5em,
        after skip=1em,
        title=\textbf{Solution},
        coltitle=uwpurple,
        fonttitle=\bfseries
    ]
}{
    \end{tcolorbox}
}

%%%%%%%%%%%%%%%%%%%%%%%%%%%%%%%%%%%%%
% Optional Source Command
%%%%%%%%%%%%%%%%%%%%%%%%%%%%%%%%%%%%%

\newcommand{\source}[1]{%
    \vspace{0.1cm}
    
    {\small\textit{Source: #1}}
}

%%%%%%%%%%%%%%%%%%%%%%%%%%%%%%%%%%%%%
% Optional Answer Boxes
%%%%%%%%%%%%%%%%%%%%%%%%%%%%%%%%%%%%%

\newcommand{\answerbox}{%
    \begin{flushright}
        \begin{tcolorbox}[
            width=0.36\textwidth,
            height=2cm,
            colback=white,
            colframe=uwpurple,
            boxrule=0.7pt,
            arc=1mm
        ]
            {\small\textbf{Final Answer}}
        \end{tcolorbox}
    \end{flushright}
}

\newcommand{\partanswerbox}[1]{%
    \begin{flushright}
        \begin{tcolorbox}[
            width=0.36\textwidth,
            height=2cm,
            colback=white,
            colframe=uwpurple,
            boxrule=0.7pt,
            arc=1mm
        ]
            {\small\textbf{Final Answer for Part #1}}
        \end{tcolorbox}
    \end{flushright}
}

%%%%%%%%%%%%%%%%%%%%%%%%%%%%%%%%%%%%%
% Mathematical Commands
%%%%%%%%%%%%%%%%%%%%%%%%%%%%%%%%%%%%%

\newcommand{\E}{\mathbb{E}}
\newcommand{\Pp}{\mathbb{P}}
\newcommand{\cC}{\mathcal{C}}

%%%%%%%%%%%%%%%%%%%%%%%%%%%%%%%%%%%%%
% Document
%%%%%%%%%%%%%%%%%%%%%%%%%%%%%%%%%%%%%

\begin{document}

%%%%%%%%%%%%%%%%%%%%%%%%%%%%%%%%%%%%%
% Title Page
%%%%%%%%%%%%%%%%%%%%%%%%%%%%%%%%%%%%%

\begin{titlepage}
    \thispagestyle{empty}

    \begin{center}

        \vspace*{1.2in}

        {\color{uwpurple}\rule{\textwidth}{1.5pt}}

        \vspace{0.45in}

        {\LARGE\bfseries\color{uwpurple}
        \documenttitle}

        \vspace{0.25in}

        {\Large\bfseries
        \course}

        \vspace{0.45in}

        {\color{uwgold}\rule{0.65\textwidth}{1.2pt}}

        \vfill

        \begin{tcolorbox}[
            width=0.72\textwidth,
            colback=softpurple,
            colframe=uwpurple,
            boxrule=0.8pt,
            arc=2mm,
            left=5mm,
            right=5mm,
            top=4mm,
            bottom=4mm
        ]
            \centering
            \large
            \textbf{Instructor}\\
            \instructor

            \vspace{0.3cm}

            \textbf{Term}\\
            \term
        \end{tcolorbox}

        \vfill

        {\large\bfseries\university}

        \vspace{0.1in}

        {\large\department}

        \vspace{0.6in}

        {\color{uwpurple}\rule{\textwidth}{1.5pt}}

    \end{center}
\end{titlepage}


























%%%%%%%%%%%%%%%%%%%%%%%%%%%%%%%%%%%%%%%%%%%%%%%%%%%%%%%%%%%%%%%%%%%%%%%%%%%%%%%
% FINAL EXAM QUESTION BANK
%%%%%%%%%%%%%%%%%%%%%%%%%%%%%%%%%%%%%%%%%%%%%%%%%%%%%%%%%%%%%%%%%%%%%%%%%%%%%%%

\newpage


%%%%%%%%%%%%%%%%%%%%%%%%%%%%%%%%%%%%%%%%%%%%%%%%%%%%%%%%%%%%%%
%%%%%%%%%%%%%%%%%%%%%%  PROBLEM 1  %%%%%%%%%%%%%%%%%%%%%%%%%%%
%%%%%%%%%%%%%%%%%%%%%%%%%%%%%%%%%%%%%%%%%%%%%%%%%%%%%%%%%%%%%%



%%%%%%%%%%%%%%%%%%%%%%%%%%%%%%%%%%%%%%%%%%%%%%%%%%%%%%%%%%%%%%
%%%%%%%%%%%%%%%%%%%%%%  MOMENTS PROBLEM  %%%%%%%%%%%%%%%%%%%%%
%%%%%%%%%%%%%%%%%%%%%%%%%%%%%%%%%%%%%%%%%%%%%%%%%%%%%%%%%%%%%%

\newpage
\begin{problem}[]

Let \(X\) be a continuous random variable with probability density function
\[
f_X(x)
=
\begin{cases}
2x, & 0<x<1,\\
0, & \text{otherwise}.
\end{cases}
\]

\begin{enumerate}[label=(\alph*)]

    \item Find the general \(k\)-th raw moment
    \[
    \E(X^k),
    \qquad k=1,2,\ldots
    \]

    \item Let \(M_X(t)\) denote the moment generating function of \(X\).
    Using your result from part (a), find
    \[
    M_X^{(k)}(0),
    \qquad k=1,2,\ldots,
    \]
    where \(M_X^{(k)}\) denotes the \(k\)-th derivative of \(M_X\).

    In particular, find
    \[
    M_X'(0),\qquad
    M_X''(0),\qquad
    M_X^{(3)}(0),\qquad
    M_X^{(4)}(0).
    \]

    \item Find
    \[
    \operatorname{Var}(X).
    \]

    \item Find the median \(m\) of \(X\).

    \item Find the skewness
    \[
    \gamma_1
    =
    \frac{\mathbb{E}[(X-\mu)^3]}{\sigma^3},
    \]
    where
    \[
    \mu=\E(X),
    \qquad
    \sigma^2=\operatorname{Var}(X).
    \]

    \item Find the kurtosis
    \[
    \gamma_2
    =
    \frac{\mathbb{E}[(X-\mu)^4]}{\sigma^4}.
    \]

\end{enumerate}

\end{problem}


\begin{solution}

\begin{enumerate}[label=(\alph*)]

\item For \(k\geq1\),
\[
\begin{aligned}
\E(X^k)
&=
\int_{-\infty}^{\infty}x^k f_X(x)\,dx\\
&=
\int_0^1 x^k(2x)\,dx\\
&=
2\int_0^1 x^{k+1}\,dx\\
&=
2\left[\frac{x^{k+2}}{k+2}\right]_0^1\\
&=
\frac{2}{k+2}.
\end{aligned}
\]

Therefore,
\[
\E(X^k)=\frac{2}{k+2}.
\]


\item Recall that, whenever the moment generating function exists in a
neighborhood of \(0\),
\[
M_X^{(k)}(0)=\E(X^k).
\]

From part (a),
\[
\E(X^k)=\frac{2}{k+2}.
\]

Therefore,
\[
M_X^{(k)}(0)
=
\frac{2}{k+2},
\qquad k=1,2,\ldots
\]

In particular,
\[
M_X'(0)=\frac23,
\qquad
M_X''(0)=\frac12,
\qquad
M_X^{(3)}(0)=\frac25,
\qquad
M_X^{(4)}(0)=\frac13.
\]


\item We have
\[
\E(X)=M_X'(0)=\frac23
\]
and
\[
\E(X^2)=M_X''(0)=\frac12.
\]

Therefore,
\[
\begin{aligned}
\operatorname{Var}(X)
&=
\E(X^2)-[\E(X)]^2\\
&=
\frac12-\left(\frac23\right)^2\\
&=
\frac12-\frac49\\
&=
\frac1{18}.
\end{aligned}
\]

Thus,
\[
\sigma^2=\frac1{18},
\qquad
\sigma=\frac{1}{3\sqrt2}.
\]


\item First calculate the CDF. For \(0\leq x\leq1\),
\[
\begin{aligned}
F_X(x)
&=
P(X\leq x)\\
&=
\int_0^x 2t\,dt\\
&=
x^2.
\end{aligned}
\]

The median \(m\) satisfies
\[
F_X(m)=\frac12.
\]

Hence,
\[
m^2=\frac12.
\]

Since \(m\in(0,1)\),
\[
m=\frac1{\sqrt2}.
\]


\item The third central moment is
\[
\mathbb{E}[(X-\mu)^3].
\]

Expanding,
\[
(X-\mu)^3
=
X^3-3\mu X^2+3\mu^2X-\mu^3.
\]

Therefore,
\[
\mathbb{E}[(X-\mu)^3]
=
\E(X^3)
-3\mu\E(X^2)
+3\mu^2\E(X)
-\mu^3.
\]

Using
\[
\mu=\frac23,
\qquad
\E(X^2)=\frac12,
\qquad
\E(X^3)=\frac25,
\]
we obtain
\[
\begin{aligned}
\mathbb{E}[(X-\mu)^3]
&=
\frac25
-
3\left(\frac23\right)\left(\frac12\right)
+
3\left(\frac23\right)^2\left(\frac23\right)
-
\left(\frac23\right)^3\\
&=
\frac25-1+\frac89-\frac8{27}\\
&=
-\frac1{135}.
\end{aligned}
\]

Also,
\[
\sigma^3
=
\left(\frac1{18}\right)^{3/2}
=
\frac{1}{54\sqrt2}.
\]

Hence,
\[
\begin{aligned}
\gamma_1
&=
\frac{\mathbb{E}[(X-\mu)^3]}{\sigma^3}\\
&=
\frac{-1/135}{1/(54\sqrt2)}\\
&=
-\frac{2\sqrt2}{5}.
\end{aligned}
\]

Therefore,
\[
\gamma_1=-\frac{2\sqrt2}{5}.
\]

Since the skewness is negative, the distribution is negatively skewed.


\item The fourth central moment is
\[
\mathbb{E}[(X-\mu)^4].
\]

Expanding,
\[
(X-\mu)^4
=
X^4
-4\mu X^3
+6\mu^2X^2
-4\mu^3X
+\mu^4.
\]

Therefore,
\[
\mathbb{E}[(X-\mu)^4]
=
\E(X^4)
-4\mu\E(X^3)
+6\mu^2\E(X^2)
-4\mu^3\E(X)
+\mu^4.
\]

Using
\[
\mu=\frac23,
\qquad
\E(X^2)=\frac12,
\qquad
\E(X^3)=\frac25,
\qquad
\E(X^4)=\frac13,
\]
we obtain
\[
\begin{aligned}
\mathbb{E}[(X-\mu)^4]
&=
\frac13
-
4\left(\frac23\right)\left(\frac25\right)
+
6\left(\frac23\right)^2\left(\frac12\right)\\
&\qquad
-
4\left(\frac23\right)^3\left(\frac23\right)
+
\left(\frac23\right)^4\\
&=
\frac13
-\frac{16}{15}
+\frac43
-\frac{64}{81}
+\frac{16}{81}\\
&=
\frac1{135}.
\end{aligned}
\]

Also,
\[
\sigma^4
=
\left(\frac1{18}\right)^2
=
\frac1{324}.
\]

Hence,
\[
\begin{aligned}
\gamma_2
&=
\frac{\mathbb{E}[(X-\mu)^4]}{\sigma^4}\\
&=
\frac{1/135}{1/324}\\
&=
\frac{12}{5}.
\end{aligned}
\]

Therefore,
\[
\gamma_2=\frac{12}{5}.
\]

\end{enumerate}

\end{solution}



%%%%%%%%%%%%%%%%%%%%%%%%%%%%%%%%%%%%%%%%%%%%%%%%%%%%%%%%%%%%%%
%%%%%%%%%%%%%%%%%%%%%%  MGF PROBLEM  %%%%%%%%%%%%%%%%%%%%%%%%%
%%%%%%%%%%%%%%%%%%%%%%%%%%%%%%%%%%%%%%%%%%%%%%%%%%%%%%%%%%%%%%

\newpage
\begin{problem}[ ]

Let \(X\) and \(Y\) be independent random variables with moment generating
functions
\[
M_X(t)=e^{2(e^t-1)}
\]
and
\[
M_Y(t)=\frac{3}{3-t},
\qquad t<3.
\]

\begin{enumerate}[label=(\alph*)]

    \item Find the moment generating function of
    \[
    X+Y.
    \]

    \item More generally, let \(a,b\in\mathbb{R}\). Find the moment
    generating function of
    \[
    Z=aX+bY.
    \]

    State the values of \(t\) for which your answer is valid.

\end{enumerate}

\end{problem}


\begin{solution}

\begin{enumerate}[label=(\alph*)]

\item Since \(X\) and \(Y\) are independent,
\[
\begin{aligned}
M_{X+Y}(t)
&=
\E\left(e^{t(X+Y)}\right)\\
&=
\E\left(e^{tX}e^{tY}\right)\\
&=
\E\left(e^{tX}\right)
\E\left(e^{tY}\right)\\
&=
M_X(t)M_Y(t).
\end{aligned}
\]

Therefore,
\[
M_{X+Y}(t)
=
e^{2(e^t-1)}
\frac{3}{3-t},
\qquad t<3.
\]


\item Let
\[
Z=aX+bY.
\]

By definition,
\[
\begin{aligned}
M_Z(t)
&=
\E\left(e^{tZ}\right)\\
&=
\E\left(e^{t(aX+bY)}\right)\\
&=
\E\left(e^{atX}e^{btY}\right).
\end{aligned}
\]

Since \(X\) and \(Y\) are independent,
\[
\begin{aligned}
M_Z(t)
&=
\E\left(e^{atX}\right)
\E\left(e^{btY}\right)\\
&=
M_X(at)M_Y(bt).
\end{aligned}
\]

Now,
\[
M_X(at)
=
e^{2(e^{at}-1)}
\]
and
\[
M_Y(bt)
=
\frac{3}{3-bt}.
\]

Hence,
\[
M_{aX+bY}(t)
=
e^{2(e^{at}-1)}
\frac{3}{3-bt}.
\]

The restriction comes from \(M_Y(bt)\), which requires
\[
bt<3.
\]

Thus,
\[
M_{aX+bY}(t)
=
\frac{3e^{2(e^{at}-1)}}{3-bt},
\qquad bt<3.
\]

Equivalently, the domain is
\[
\begin{cases}
t<3/b, & b>0,\\
t>3/b, & b<0,\\
t\in\mathbb{R}, & b=0.
\end{cases}
\]

\end{enumerate}

\end{solution}





























\newpage
\begin{problem}

Let the joint probability mass function of discrete random variables \(X\) and
\(Y\) be given by
\[
p(x,y)
=
\begin{cases}
c(x+y),
&
x=1,2,3,\quad y=1,2,\\[1mm]
0,
&
\text{otherwise}.
\end{cases}
\]

Determine

\begin{enumerate}[label=(\alph*)]
    \item the value of the constant \(c\);
    \item the marginal probability mass functions of \(X\) and \(Y\);
    \item \(P(X\geq2\mid Y=1)\);
    \item \(E(X)\) and \(E(Y)\).
\end{enumerate}

\end{problem}

\begin{solution}

\begin{enumerate}[label=(\alph*)]

\item Since the total probability must equal \(1\),
\[
\sum_{x=1}^{3}\sum_{y=1}^{2}c(x+y)=1.
\]

Thus,
\[
c[(2+3)+(3+4)+(4+5)]=1.
\]

Therefore,
\[
21c=1,
\]
and hence
\[
c=\frac1{21}.
\]

\item For \(x=1,2,3\),
\[
\begin{aligned}
p_X(x)
&=
\sum_{y=1}^{2}p(x,y)\\
&=
\frac1{21}\left[(x+1)+(x+2)\right]\\
&=
\frac{2x+3}{21}.
\end{aligned}
\]

Thus,
\[
p_X(x)
=
\begin{cases}
\dfrac{2x+3}{21}, & x=1,2,3,\\[1mm]
0, & \text{otherwise}.
\end{cases}
\]

For \(y=1,2\),
\[
\begin{aligned}
p_Y(y)
&=
\sum_{x=1}^{3}p(x,y)\\
&=
\frac1{21}\left[(1+y)+(2+y)+(3+y)\right]\\
&=
\frac{6+3y}{21}.
\end{aligned}
\]

Therefore,
\[
p_Y(y)
=
\begin{cases}
\dfrac37, & y=1,\\[1mm]
\dfrac47, & y=2,\\[1mm]
0, & \text{otherwise}.
\end{cases}
\]

\item
\[
P(X\geq2\mid Y=1)
=
\frac{P(X\geq2,Y=1)}{P(Y=1)}.
\]

The numerator is
\[
\begin{aligned}
P(X\geq2,Y=1)
&=
p(2,1)+p(3,1)\\
&=
\frac3{21}+\frac4{21}\\
&=
\frac7{21}.
\end{aligned}
\]

Also,
\[
P(Y=1)=\frac37=\frac9{21}.
\]

Therefore,
\[
P(X\geq2\mid Y=1)
=
\frac{7/21}{9/21}
=
\frac79.
\]

\item
\[
\begin{aligned}
E(X)
&=
\sum_{x=1}^{3}x\,p_X(x)\\
&=
1\left(\frac5{21}\right)
+
2\left(\frac7{21}\right)
+
3\left(\frac9{21}\right)\\
&=
\frac{46}{21}.
\end{aligned}
\]

Similarly,
\[
E(Y)
=
1\left(\frac37\right)
+
2\left(\frac47\right)
=
\frac{11}{7}.
\]

Thus,
\[
E(X)=\frac{46}{21},
\qquad
E(Y)=\frac{11}{7}.
\]

\end{enumerate}

\end{solution}


%%%%%%%%%%%%%%%%%%%%%%%%%%%%%%%%%%%%%%%%%%%%%%%%%%%%%%%%%%%%%%
%%%%%%%%%%%%%%%%%%%%%%  PROBLEM 3  %%%%%%%%%%%%%%%%%%%%%%%%%%%
%%%%%%%%%%%%%%%%%%%%%%%%%%%%%%%%%%%%%%%%%%%%%%%%%%%%%%%%%%%%%%

\newpage
\begin{problem}

Let the joint probability density function of random variables \(X\) and \(Y\)
be given by
\[
f(x,y)
=
\begin{cases}
8xy,
&
0\leq y\leq x\leq1,\\
0,
&
\text{elsewhere}.
\end{cases}
\]

\begin{enumerate}[label=(\alph*)]
    \item Calculate the marginal probability density functions of \(X\) and
    \(Y\).
    \item Calculate \(E(X)\) and \(E(Y)\).
\end{enumerate}

\end{problem}

\begin{solution}

The support of the joint density is
\[
0\leq y\leq x\leq1.
\]

\begin{enumerate}[label=(\alph*)]

\item For a fixed \(x\in[0,1]\), the possible values of \(y\) satisfy
\[
0\leq y\leq x.
\]

Therefore,
\[
\begin{aligned}
f_X(x)
&=
\int_0^x 8xy\,dy\\
&=
8x\left[\frac{y^2}{2}\right]_0^x\\
&=
4x^3.
\end{aligned}
\]

Hence,
\[
f_X(x)
=
\begin{cases}
4x^3, & 0\leq x\leq1,\\
0, & \text{otherwise}.
\end{cases}
\]

For a fixed \(y\in[0,1]\), the support requires
\[
y\leq x\leq1.
\]

Thus,
\[
\begin{aligned}
f_Y(y)
&=
\int_y^1 8xy\,dx\\
&=
8y\left[\frac{x^2}{2}\right]_y^1\\
&=
4y(1-y^2).
\end{aligned}
\]

Therefore,
\[
f_Y(y)
=
\begin{cases}
4y(1-y^2), & 0\leq y\leq1,\\
0, & \text{otherwise}.
\end{cases}
\]

\item Using the marginal density of \(X\),
\[
\begin{aligned}
E(X)
&=
\int_0^1 x f_X(x)\,dx\\
&=
\int_0^1 4x^4\,dx\\
&=
\frac45.
\end{aligned}
\]

Similarly,
\[
\begin{aligned}
E(Y)
&=
\int_0^1 y f_Y(y)\,dy\\
&=
4\int_0^1 y^2(1-y^2)\,dy\\
&=
4\left(\frac13-\frac15\right)\\
&=
\frac8{15}.
\end{aligned}
\]

Thus,
\[
E(X)=\frac45,
\qquad
E(Y)=\frac8{15}.
\]

\end{enumerate}

\end{solution}


%%%%%%%%%%%%%%%%%%%%%%%%%%%%%%%%%%%%%%%%%%%%%%%%%%%%%%%%%%%%%%
%%%%%%%%%%%%%%%%%%%%%%  PROBLEM 4  %%%%%%%%%%%%%%%%%%%%%%%%%%%
%%%%%%%%%%%%%%%%%%%%%%%%%%%%%%%%%%%%%%%%%%%%%%%%%%%%%%%%%%%%%%

\newpage
\begin{problem}

Let the joint probability density function of random variables \(X\) and \(Y\)
be given by
\[
f(x,y)
=
\begin{cases}
\dfrac12 ye^{-x},
&
x>0,\quad 0<y<2,\\[2mm]
0,
&
\text{elsewhere}.
\end{cases}
\]

Find the marginal probability density functions of \(X\) and \(Y\).

\end{problem}

\begin{solution}

For \(x>0\),
\[
\begin{aligned}
f_X(x)
&=
\int_0^2 \frac12 ye^{-x}\,dy\\
&=
\frac12e^{-x}
\left[\frac{y^2}{2}\right]_0^2\\
&=
e^{-x}.
\end{aligned}
\]

Thus,
\[
f_X(x)
=
\begin{cases}
e^{-x}, & x>0,\\
0, & \text{otherwise}.
\end{cases}
\]

For \(0<y<2\),
\[
\begin{aligned}
f_Y(y)
&=
\int_0^\infty \frac12 ye^{-x}\,dx\\
&=
\frac y2\int_0^\infty e^{-x}\,dx\\
&=
\frac y2.
\end{aligned}
\]

Therefore,
\[
f_Y(y)
=
\begin{cases}
\dfrac y2, & 0<y<2,\\[1mm]
0, & \text{otherwise}.
\end{cases}
\]

\end{solution}


%%%%%%%%%%%%%%%%%%%%%%%%%%%%%%%%%%%%%%%%%%%%%%%%%%%%%%%%%%%%%%
%%%%%%%%%%%%%%%%%%%%%%  PROBLEM 5  %%%%%%%%%%%%%%%%%%%%%%%%%%%
%%%%%%%%%%%%%%%%%%%%%%%%%%%%%%%%%%%%%%%%%%%%%%%%%%%%%%%%%%%%%%

\newpage
\begin{problem}

Let \(X\) and \(Y\) have the joint probability density function
\[
f(x,y)
=
\begin{cases}
1,
&
0\leq x\leq1,\quad 0\leq y\leq1,\\
0,
&
\text{elsewhere}.
\end{cases}
\]

Calculate
\[
P\left(X+Y\leq\frac12\right),
\qquad
P\left(X-Y\leq\frac12\right),
\]
\[
P\left(XY\leq\frac14\right),
\qquad
P(X^2+Y^2\leq1).
\]

\end{problem}

\begin{solution}

Since the density is equal to \(1\) on the unit square, each desired
probability is equal to the area of the corresponding region inside
\([0,1]^2\).

First,
\[
X+Y\leq\frac12
\]
describes a right triangle with legs of length \(1/2\). Hence,
\[
P\left(X+Y\leq\frac12\right)
=
\frac12\left(\frac12\right)\left(\frac12\right)
=
\frac18.
\]

Next,
\[
X-Y\leq\frac12.
\]

It is easier to subtract the region
\[
X-Y>\frac12.
\]

This is a right triangle in the lower-right portion of the unit square with
legs of length \(1/2\). Therefore,
\[
P\left(X-Y>\frac12\right)=\frac18,
\]
and hence
\[
P\left(X-Y\leq\frac12\right)
=
1-\frac18
=
\frac78.
\]

For the third probability,
\[
XY\leq\frac14.
\]

If \(0\leq x\leq1/4\), every \(0\leq y\leq1\) satisfies the inequality.

If \(1/4<x\leq1\), then
\[
y\leq\frac{1}{4x}.
\]

Thus,
\[
\begin{aligned}
P\left(XY\leq\frac14\right)
&=
\int_0^{1/4}\int_0^1dy\,dx
+
\int_{1/4}^{1}\int_0^{1/(4x)}dy\,dx\\
&=
\frac14
+
\frac14\int_{1/4}^{1}\frac1x\,dx\\
&=
\frac14+\frac14\ln4.
\end{aligned}
\]

Finally,
\[
X^2+Y^2\leq1
\]
describes the quarter of the unit disk lying in the first quadrant.
Therefore,
\[
P(X^2+Y^2\leq1)
=
\frac{\pi}{4}.
\]

Hence,
\[
P\left(X+Y\leq\frac12\right)=\frac18,
\]
\[
P\left(X-Y\leq\frac12\right)=\frac78,
\]
\[
P\left(XY\leq\frac14\right)
=
\frac14+\frac14\ln4,
\]
and
\[
P(X^2+Y^2\leq1)=\frac{\pi}{4}.
\]

\end{solution}


%%%%%%%%%%%%%%%%%%%%%%%%%%%%%%%%%%%%%%%%%%%%%%%%%%%%%%%%%%%%%%
%%%%%%%%%%%%%%%%%%%%%%  PROBLEM 6  %%%%%%%%%%%%%%%%%%%%%%%%%%%
%%%%%%%%%%%%%%%%%%%%%%%%%%%%%%%%%%%%%%%%%%%%%%%%%%%%%%%%%%%%%%

\newpage
\begin{problem}

Suppose that \(h\) is the probability density function of a continuous random
variable. Let the joint probability density function of \(X\) and \(Y\) be
given by
\[
f(x,y)=h(x)h(y),
\qquad
x,y\in\mathbb R.
\]

Prove that
\[
P(X\geq Y)=\frac12.
\]

\end{problem}

\begin{solution}

Since
\[
f(x,y)=h(x)h(y),
\]
\(X\) and \(Y\) are independent and identically distributed continuous random
variables.

Because \(X\) and \(Y\) have the same distribution, symmetry gives
\[
P(X>Y)=P(Y>X).
\]

Moreover, since \(X\) and \(Y\) are continuous,
\[
P(X=Y)=0.
\]

The events
\[
\{X>Y\},
\qquad
\{Y>X\},
\qquad
\{X=Y\}
\]
partition the sample space. Therefore,
\[
P(X>Y)+P(Y>X)+P(X=Y)=1.
\]

Hence,
\[
2P(X>Y)=1,
\]
so
\[
P(X>Y)=\frac12.
\]

Finally,
\[
P(X\geq Y)
=
P(X>Y)+P(X=Y)
=
\frac12.
\]

\end{solution}


%%%%%%%%%%%%%%%%%%%%%%%%%%%%%%%%%%%%%%%%%%%%%%%%%%%%%%%%%%%%%%
%%%%%%%%%%%%%%%%%%%%%%  PROBLEM 7  %%%%%%%%%%%%%%%%%%%%%%%%%%%
%%%%%%%%%%%%%%%%%%%%%%%%%%%%%%%%%%%%%%%%%%%%%%%%%%%%%%%%%%%%%%

\newpage
\begin{problem}

Let \(X\) and \(Y\) be continuous random variables with joint probability
density function \(f(x,y)\). Define
\[
Z=\frac{Y}{X},
\qquad X\neq0.
\]

Prove that the probability density function of \(Z\) is
\[
f_Z(z)
=
\int_{-\infty}^{\infty}
|x|f(x,xz)\,dx.
\]

\end{problem}

\begin{solution}

Introduce the transformation
\[
Z=\frac{Y}{X},
\qquad
W=X.
\]

Then
\[
z=\frac{y}{x},
\qquad
w=x.
\]

Solving for \(x\) and \(y\) in terms of \(z\) and \(w\),
\[
x=w,
\qquad
y=wz.
\]

The Jacobian matrix of the inverse transformation is
\[
\frac{\partial(x,y)}{\partial(z,w)}
=
\begin{pmatrix}
0 & 1\\
w & z
\end{pmatrix}.
\]

Therefore,
\[
\det
\frac{\partial(x,y)}{\partial(z,w)}
=
-w,
\]
and hence
\[
\left|
\det
\frac{\partial(x,y)}{\partial(z,w)}
\right|
=
|w|.
\]

By the change-of-variables formula,
\[
f_{Z,W}(z,w)
=
f(w,wz)|w|.
\]

To obtain the marginal density of \(Z\), integrate out \(W\):
\[
\begin{aligned}
f_Z(z)
&=
\int_{-\infty}^{\infty}
f_{Z,W}(z,w)\,dw\\
&=
\int_{-\infty}^{\infty}
|w|f(w,wz)\,dw.
\end{aligned}
\]

Renaming the variable of integration from \(w\) to \(x\),
\[
f_Z(z)
=
\int_{-\infty}^{\infty}
|x|f(x,xz)\,dx.
\]

\end{solution}


%%%%%%%%%%%%%%%%%%%%%%%%%%%%%%%%%%%%%%%%%%%%%%%%%%%%%%%%%%%%%%
%%%%%%%%%%%%%%%%%%%%%%  PROBLEM 8  %%%%%%%%%%%%%%%%%%%%%%%%%%%
%%%%%%%%%%%%%%%%%%%%%%%%%%%%%%%%%%%%%%%%%%%%%%%%%%%%%%%%%%%%%%

\newpage
\begin{problem}

Prove that random variables \(X\) and \(Y\) with joint probability density
function
\[
f(x,y)
=
\begin{cases}
8xy,
&
0\leq x\leq y\leq1,\\
0,
&
\text{otherwise},
\end{cases}
\]
are not independent.

\end{problem}

\begin{solution}

We first calculate the marginal densities.

For \(0\leq x\leq1\),
\[
\begin{aligned}
f_X(x)
&=
\int_x^1 8xy\,dy\\
&=
8x\left[\frac{y^2}{2}\right]_x^1\\
&=
4x(1-x^2).
\end{aligned}
\]

Thus,
\[
f_X(x)
=
\begin{cases}
4x(1-x^2), & 0\leq x\leq1,\\
0, & \text{otherwise}.
\end{cases}
\]

For \(0\leq y\leq1\),
\[
\begin{aligned}
f_Y(y)
&=
\int_0^y8xy\,dx\\
&=
8y\left[\frac{x^2}{2}\right]_0^y\\
&=
4y^3.
\end{aligned}
\]

Hence,
\[
f_Y(y)
=
\begin{cases}
4y^3, & 0\leq y\leq1,\\
0, & \text{otherwise}.
\end{cases}
\]

If \(X\) and \(Y\) were independent, we would have
\[
f(x,y)=f_X(x)f_Y(y)
\]
for all \(x,y\).

However,
\[
f_X(x)f_Y(y)
=
16xy^3(1-x^2),
\]
which is not equal to
\[
f(x,y)=8xy
\]
on the support.

Therefore, \(X\) and \(Y\) are not independent.

There is also a geometric reason for the dependence. The support
\[
0\leq x\leq y\leq1
\]
is triangular rather than a Cartesian product of one-dimensional supports.
In particular, knowing the value of \(Y\) restricts the possible values of
\(X\).

\end{solution}


%%%%%%%%%%%%%%%%%%%%%%%%%%%%%%%%%%%%%%%%%%%%%%%%%%%%%%%%%%%%%%
%%%%%%%%%%%%%%%%%%%%%%  PROBLEM 9  %%%%%%%%%%%%%%%%%%%%%%%%%%%
%%%%%%%%%%%%%%%%%%%%%%%%%%%%%%%%%%%%%%%%%%%%%%%%%%%%%%%%%%%%%%

\newpage
\begin{problem}

Let the joint probability mass function of random variables \(X\) and \(Y\)
be given by
\[
p(x,y)
=
\begin{cases}
\dfrac17x^2y,
&
(x,y)=(1,1),(1,2),(2,1),\\[2mm]
0,
&
\text{elsewhere}.
\end{cases}
\]

Are \(X\) and \(Y\) independent? Why or why not?

\end{problem}

\begin{solution}

The nonzero joint probabilities are
\[
p(1,1)=\frac17,
\qquad
p(1,2)=\frac27,
\qquad
p(2,1)=\frac47.
\]

Therefore,
\[
P(X=1)
=
\frac17+\frac27
=
\frac37,
\]
and
\[
P(X=2)=\frac47.
\]

Similarly,
\[
P(Y=1)
=
\frac17+\frac47
=
\frac57,
\]
and
\[
P(Y=2)=\frac27.
\]

If \(X\) and \(Y\) were independent, then
\[
P(X=2,Y=2)
=
P(X=2)P(Y=2).
\]

However,
\[
P(X=2,Y=2)=0,
\]
whereas
\[
P(X=2)P(Y=2)
=
\frac47\frac27
=
\frac8{49}>0.
\]

Therefore,
\[
P(X=2,Y=2)
\neq
P(X=2)P(Y=2),
\]
so \(X\) and \(Y\) are not independent.

\end{solution}


%%%%%%%%%%%%%%%%%%%%%%%%%%%%%%%%%%%%%%%%%%%%%%%%%%%%%%%%%%%%%%
%%%%%%%%%%%%%%%%%%%%%%  PROBLEM 10  %%%%%%%%%%%%%%%%%%%%%%%%%%
%%%%%%%%%%%%%%%%%%%%%%%%%%%%%%%%%%%%%%%%%%%%%%%%%%%%%%%%%%%%%%

\newpage
\begin{problem}

Let the joint probability density function of random variables \(X\) and \(Y\)
be given by
\[
f(x,y)
=
\begin{cases}
x^2e^{-x(y+1)},
&
x\geq0,\quad y\geq0,\\
0,
&
\text{elsewhere}.
\end{cases}
\]

Are \(X\) and \(Y\) independent? Why or why not?

\end{problem}

\begin{solution}

We calculate the marginal densities.

For \(x>0\),
\[
\begin{aligned}
f_X(x)
&=
\int_0^\infty x^2e^{-x(y+1)}\,dy\\
&=
x^2e^{-x}
\int_0^\infty e^{-xy}\,dy\\
&=
x^2e^{-x}\left(\frac1x\right)\\
&=
xe^{-x}.
\end{aligned}
\]

Thus,
\[
f_X(x)
=
\begin{cases}
xe^{-x}, & x>0,\\
0, & \text{otherwise}.
\end{cases}
\]

For \(y\geq0\),
\[
\begin{aligned}
f_Y(y)
&=
\int_0^\infty x^2e^{-(y+1)x}\,dx.
\end{aligned}
\]

Using
\[
\int_0^\infty x^2e^{-ax}\,dx
=
\frac{2}{a^3},
\qquad a>0,
\]
we obtain
\[
f_Y(y)
=
\frac{2}{(y+1)^3},
\qquad y\geq0.
\]

Therefore,
\[
f_X(x)f_Y(y)
=
\frac{2xe^{-x}}{(y+1)^3}.
\]

But the joint density is
\[
f(x,y)
=
x^2e^{-x(y+1)}.
\]

These are not equal in general. Hence,
\[
f(x,y)\neq f_X(x)f_Y(y),
\]
and therefore \(X\) and \(Y\) are not independent.

\end{solution}


%%%%%%%%%%%%%%%%%%%%%%%%%%%%%%%%%%%%%%%%%%%%%%%%%%%%%%%%%%%%%%
%%%%%%%%%%%%%%%%%%%%%%  PROBLEM 11  %%%%%%%%%%%%%%%%%%%%%%%%%%
%%%%%%%%%%%%%%%%%%%%%%%%%%%%%%%%%%%%%%%%%%%%%%%%%%%%%%%%%%%%%%

\newpage
\begin{problem}

Suppose that \(X\) and \(Y\) are independent, identically distributed
exponential random variables with mean \(1/\lambda\). Prove that
\[
\frac{X}{X+Y}
\]
is uniformly distributed on \((0,1)\).

\end{problem}

\begin{solution}

Since
\[
X,Y\overset{\text{i.i.d.}}{\sim}\operatorname{Exp}(\lambda),
\]
their joint density is
\[
f_{X,Y}(x,y)
=
\lambda^2e^{-\lambda(x+y)},
\qquad x>0,\ y>0.
\]

Define
\[
U=\frac{X}{X+Y},
\qquad
V=X+Y.
\]

Then
\[
X=UV
\]
and
\[
Y=(1-U)V.
\]

Since \(X>0\) and \(Y>0\), the support of \((U,V)\) is
\[
0<u<1,
\qquad
v>0.
\]

The Jacobian matrix is
\[
\frac{\partial(x,y)}{\partial(u,v)}
=
\begin{pmatrix}
v & u\\
-v & 1-u
\end{pmatrix}.
\]

Therefore,
\[
\det
\frac{\partial(x,y)}{\partial(u,v)}
=
v(1-u)+uv
=
v.
\]

Hence,
\[
\left|
\det
\frac{\partial(x,y)}{\partial(u,v)}
\right|
=
v.
\]

The joint density of \(U\) and \(V\) is therefore
\[
\begin{aligned}
f_{U,V}(u,v)
&=
f_{X,Y}(uv,(1-u)v)v\\
&=
\lambda^2e^{-\lambda v}v,
\end{aligned}
\]
for
\[
0<u<1,\qquad v>0.
\]

Now integrate out \(V\):
\[
\begin{aligned}
f_U(u)
&=
\int_0^\infty
\lambda^2ve^{-\lambda v}\,dv\\
&=
1,
\qquad 0<u<1.
\end{aligned}
\]

Thus,
\[
f_U(u)
=
\begin{cases}
1, & 0<u<1,\\
0, & \text{otherwise}.
\end{cases}
\]

Therefore,
\[
\frac{X}{X+Y}
\sim
\operatorname{Unif}(0,1).
\]

\end{solution}


%%%%%%%%%%%%%%%%%%%%%%%%%%%%%%%%%%%%%%%%%%%%%%%%%%%%%%%%%%%%%%
%%%%%%%%%%%%%%%%%%%%%%  PROBLEM 12  %%%%%%%%%%%%%%%%%%%%%%%%%%
%%%%%%%%%%%%%%%%%%%%%%%%%%%%%%%%%%%%%%%%%%%%%%%%%%%%%%%%%%%%%%

\newpage
\begin{problem}

Let the joint probability mass function of discrete random variables \(X\) and
\(Y\) be given by
\[
p(x,y)
=
\begin{cases}
\dfrac1{25}(x^2+y^2),
&
x=1,2,\quad y=0,1,2,\\[2mm]
0,
&
\text{otherwise}.
\end{cases}
\]

Find

\begin{enumerate}[label=(\alph*)]
    \item \(p_{X\mid Y}(x\mid y)\);
    \item \(P(X=2\mid Y=1)\);
    \item \(E(X\mid Y=1)\).
\end{enumerate}

\end{problem}

\begin{solution}

First find the marginal PMF of \(Y\). For \(y=0,1,2\),
\[
\begin{aligned}
p_Y(y)
&=
\sum_{x=1}^{2}p(x,y)\\
&=
\frac1{25}
\left[
(1+y^2)+(4+y^2)
\right]\\
&=
\frac{5+2y^2}{25}.
\end{aligned}
\]

\begin{enumerate}[label=(\alph*)]

\item By definition,
\[
p_{X\mid Y}(x\mid y)
=
\frac{p(x,y)}{p_Y(y)}.
\]

Therefore,
\[
p_{X\mid Y}(x\mid y)
=
\frac{x^2+y^2}{5+2y^2},
\qquad
x=1,2,\quad y=0,1,2.
\]

\item When \(Y=1\),
\[
P(X=2\mid Y=1)
=
\frac{2^2+1^2}{5+2(1)^2}
=
\frac57.
\]

\item When \(Y=1\),
\[
P(X=1\mid Y=1)
=
\frac{1^2+1^2}{7}
=
\frac27,
\]
and
\[
P(X=2\mid Y=1)
=
\frac57.
\]

Hence,
\[
\begin{aligned}
E(X\mid Y=1)
&=
1\left(\frac27\right)
+
2\left(\frac57\right)\\
&=
\frac{12}{7}.
\end{aligned}
\]

\end{enumerate}

\end{solution}


%%%%%%%%%%%%%%%%%%%%%%%%%%%%%%%%%%%%%%%%%%%%%%%%%%%%%%%%%%%%%%
%%%%%%%%%%%%%%%%%%%%%%  PROBLEM 13  %%%%%%%%%%%%%%%%%%%%%%%%%%
%%%%%%%%%%%%%%%%%%%%%%%%%%%%%%%%%%%%%%%%%%%%%%%%%%%%%%%%%%%%%%

\newpage
\begin{problem}

Let the conditional probability density function of \(X\) given \(Y=y\) be
\[
f_{X\mid Y}(x\mid y)
=
\frac{3(x^2+y^2)}{3y^2+1},
\qquad
0<x<1,\quad 0<y<1.
\]

Find
\[
P\left(
\frac14<X<\frac12
\,\middle|\,
Y=\frac34
\right).
\]

\end{problem}

\begin{solution}

Substituting
\[
y=\frac34
\]
into the conditional density gives
\[
f_{X\mid Y}\left(x\mid\frac34\right)
=
\frac{
3\left(x^2+\frac9{16}\right)
}{
3\left(\frac9{16}\right)+1
}.
\]

Since
\[
3\left(\frac9{16}\right)+1
=
\frac{43}{16},
\]
we obtain
\[
f_{X\mid Y}\left(x\mid\frac34\right)
=
\frac{48}{43}
\left(x^2+\frac9{16}\right).
\]

Therefore,
\[
\begin{aligned}
P\left(
\frac14<X<\frac12
\,\middle|\,
Y=\frac34
\right)
&=
\int_{1/4}^{1/2}
f_{X\mid Y}\left(x\mid\frac34\right)\,dx\\
&=
\frac{48}{43}
\int_{1/4}^{1/2}
\left(x^2+\frac9{16}\right)\,dx\\
&=
\frac{48}{43}
\left[
\frac{x^3}{3}
+
\frac9{16}x
\right]_{1/4}^{1/2}\\
&=
\frac{17}{86}.
\end{aligned}
\]

Thus,
\[
P\left(
\frac14<X<\frac12
\,\middle|\,
Y=\frac34
\right)
=
\frac{17}{86}.
\]

\end{solution}


%%%%%%%%%%%%%%%%%%%%%%%%%%%%%%%%%%%%%%%%%%%%%%%%%%%%%%%%%%%%%%
%%%%%%%%%%%%%%%%%%%%%%  PROBLEM 14  %%%%%%%%%%%%%%%%%%%%%%%%%%
%%%%%%%%%%%%%%%%%%%%%%%%%%%%%%%%%%%%%%%%%%%%%%%%%%%%%%%%%%%%%%

\newpage
\begin{problem}

Let \(X\) and \(Y\) be continuous random variables with joint probability
density function
\[
f(x,y)
=
\begin{cases}
x+y,
&
0\leq x\leq1,\quad 0\leq y\leq1,\\
0,
&
\text{elsewhere}.
\end{cases}
\]

Calculate \(f_{X\mid Y}(x\mid y)\).

\end{problem}

\begin{solution}

First calculate the marginal density of \(Y\):
\[
\begin{aligned}
f_Y(y)
&=
\int_0^1(x+y)\,dx\\
&=
\left[\frac{x^2}{2}+xy\right]_0^1\\
&=
\frac12+y,
\end{aligned}
\]
for
\[
0\leq y\leq1.
\]

Therefore,
\[
\begin{aligned}
f_{X\mid Y}(x\mid y)
&=
\frac{f(x,y)}{f_Y(y)}\\
&=
\frac{x+y}{y+\frac12},
\end{aligned}
\]
for
\[
0\leq x\leq1,
\qquad
0\leq y\leq1.
\]

Thus,
\[
f_{X\mid Y}(x\mid y)
=
\begin{cases}
\dfrac{x+y}{y+\frac12},
&
0\leq x\leq1,\quad 0\leq y\leq1,\\[2mm]
0,
&
\text{otherwise}.
\end{cases}
\]

\end{solution}


%%%%%%%%%%%%%%%%%%%%%%%%%%%%%%%%%%%%%%%%%%%%%%%%%%%%%%%%%%%%%%
%%%%%%%%%%%%%%%%%%%%%%  PROBLEM 15  %%%%%%%%%%%%%%%%%%%%%%%%%%
%%%%%%%%%%%%%%%%%%%%%%%%%%%%%%%%%%%%%%%%%%%%%%%%%%%%%%%%%%%%%%

\newpage
\begin{problem}

The joint probability density function of \(X\) and \(Y\) is given by
\[
f(x,y)
=
\begin{cases}
ce^{-x},
&
x\geq0,\quad |y|<x,\\
0,
&
\text{otherwise}.
\end{cases}
\]

\begin{enumerate}[label=(\alph*)]
    \item Determine the constant \(c\).
    \item Find \(f_{X\mid Y}(x\mid y)\) and \(f_{Y\mid X}(y\mid x)\).
    \item Calculate \(E(Y\mid X=x)\) and
    \(\operatorname{Var}(Y\mid X=x)\).
\end{enumerate}

\end{problem}

\begin{solution}

The support is
\[
x\geq0,
\qquad
-x<y<x.
\]

\begin{enumerate}[label=(\alph*)]

\item The density must integrate to \(1\):
\[
\int_0^\infty
\int_{-x}^{x}
ce^{-x}\,dy\,dx
=
1.
\]

Therefore,
\[
\begin{aligned}
1
&=
c\int_0^\infty
2xe^{-x}\,dx\\
&=
2c.
\end{aligned}
\]

Hence,
\[
c=\frac12.
\]

Thus,
\[
f(x,y)
=
\frac12e^{-x}
\]
on the support.

\item First,
\[
\begin{aligned}
f_X(x)
&=
\int_{-x}^{x}\frac12e^{-x}\,dy\\
&=
xe^{-x},
\qquad x>0.
\end{aligned}
\]

Therefore,
\[
\begin{aligned}
f_{Y\mid X}(y\mid x)
&=
\frac{f(x,y)}{f_X(x)}\\
&=
\frac{\frac12e^{-x}}{xe^{-x}}\\
&=
\frac1{2x},
\end{aligned}
\]
for
\[
-x<y<x.
\]

Thus,
\[
f_{Y\mid X}(y\mid x)
=
\begin{cases}
\dfrac1{2x}, & -x<y<x,\\[1mm]
0, & \text{otherwise}.
\end{cases}
\]

Hence, conditional on \(X=x\),
\[
Y\mid X=x
\sim
\operatorname{Unif}(-x,x).
\]

Next, for a fixed \(y\), the condition
\[
|y|<x
\]
implies
\[
x>|y|.
\]

Therefore,
\[
\begin{aligned}
f_Y(y)
&=
\int_{|y|}^{\infty}
\frac12e^{-x}\,dx\\
&=
\frac12e^{-|y|}.
\end{aligned}
\]

Hence,
\[
\begin{aligned}
f_{X\mid Y}(x\mid y)
&=
\frac{f(x,y)}{f_Y(y)}\\
&=
\frac{\frac12e^{-x}}
{\frac12e^{-|y|}}\\
&=
e^{-(x-|y|)},
\end{aligned}
\]
for
\[
x>|y|.
\]

Thus,
\[
f_{X\mid Y}(x\mid y)
=
\begin{cases}
e^{-(x-|y|)}, & x>|y|,\\
0, & \text{otherwise}.
\end{cases}
\]

\item Since
\[
Y\mid X=x
\sim
\operatorname{Unif}(-x,x),
\]
symmetry gives
\[
E(Y\mid X=x)=0.
\]

The variance of a uniform random variable on \((a,b)\) is
\[
\frac{(b-a)^2}{12}.
\]

Therefore,
\[
\begin{aligned}
\operatorname{Var}(Y\mid X=x)
&=
\frac{(x-(-x))^2}{12}\\
&=
\frac{4x^2}{12}\\
&=
\frac{x^2}{3}.
\end{aligned}
\]

 

\end{enumerate}

\end{solution}


%%%%%%%%%%%%%%%%%%%%%%%%%%%%%%%%%%%%%%%%%%%%%%%%%%%%%%%%%%%%%%
%%%%%%%%%%%%%%%%%%%%%%  PROBLEM 16  %%%%%%%%%%%%%%%%%%%%%%%%%%
%%%%%%%%%%%%%%%%%%%%%%%%%%%%%%%%%%%%%%%%%%%%%%%%%%%%%%%%%%%%%%

\newpage
\begin{problem}

Let \(X\) and \(Y\) be independent random variables, each uniformly
distributed on \((0,1)\). Find the joint probability density function of
\[
U=-2\ln X
\qquad\text{and}\qquad
V=-2\ln Y.
\]

\end{problem}

\begin{solution}

Since
\[
X,Y\overset{\text{i.i.d.}}{\sim}\operatorname{Unif}(0,1),
\]
their joint density is
\[
f_{X,Y}(x,y)=1,
\qquad
0<x<1,\quad0<y<1.
\]

The transformation is
\[
u=-2\ln x,
\qquad
v=-2\ln y.
\]

Solving for \(x\) and \(y\),
\[
x=e^{-u/2},
\qquad
y=e^{-v/2}.
\]

Since \(0<x,y<1\), we have
\[
u>0,
\qquad
v>0.
\]

The Jacobian matrix of the inverse transformation is
\[
\frac{\partial(x,y)}{\partial(u,v)}
=
\begin{pmatrix}
-\dfrac12e^{-u/2} & 0\\[2mm]
0 & -\dfrac12e^{-v/2}
\end{pmatrix}.
\]

Therefore,
\[
\left|
\det
\frac{\partial(x,y)}{\partial(u,v)}
\right|
=
\frac14e^{-(u+v)/2}.
\]

Hence,
\[
f_{U,V}(u,v)
=
f_{X,Y}
\left(
e^{-u/2},e^{-v/2}
\right)
\frac14e^{-(u+v)/2}.
\]

Since the original joint density equals \(1\) on its support,
\[
f_{U,V}(u,v)
=
\begin{cases}
\dfrac14e^{-(u+v)/2},
&
u>0,\quad v>0,\\[2mm]
0,
&
\text{otherwise}.
\end{cases}
\]

\end{solution}


%%%%%%%%%%%%%%%%%%%%%%%%%%%%%%%%%%%%%%%%%%%%%%%%%%%%%%%%%%%%%%
%%%%%%%%%%%%%%%%%%%%%%  PROBLEM 17  %%%%%%%%%%%%%%%%%%%%%%%%%%
%%%%%%%%%%%%%%%%%%%%%%%%%%%%%%%%%%%%%%%%%%%%%%%%%%%%%%%%%%%%%%

\newpage
\begin{problem}

Let \(X\) and \(Y\) be two positive independent continuous random variables
with probability density functions \(f_1(x)\) and \(f_2(y)\), respectively.
Find the probability density function of
\[
U=\frac{X}{Y}.
\]

\textit{Hint:} Let \(V=X\); find the joint probability density function of
\(U\) and \(V\), and then calculate the marginal probability density function
of \(U\).

\end{problem}

\begin{solution}

Define
\[
U=\frac{X}{Y},
\qquad
V=X.
\]

Then
\[
u=\frac{x}{y},
\qquad
v=x.
\]

Solving for \(x\) and \(y\),
\[
x=v,
\qquad
y=\frac vu.
\]

Since \(X\) and \(Y\) are positive,
\[
u>0,
\qquad
v>0.
\]

The Jacobian matrix is
\[
\frac{\partial(x,y)}{\partial(u,v)}
=
\begin{pmatrix}
0 & 1\\[2mm]
-\dfrac{v}{u^2} & \dfrac1u
\end{pmatrix}.
\]

Therefore,
\[
\left|
\det
\frac{\partial(x,y)}{\partial(u,v)}
\right|
=
\frac{v}{u^2}.
\]

Since \(X\) and \(Y\) are independent,
\[
f_{X,Y}(x,y)
=
f_1(x)f_2(y).
\]

Thus,
\[
f_{U,V}(u,v)
=
f_1(v)
f_2\left(\frac vu\right)
\frac{v}{u^2},
\qquad
u>0,\quad v>0.
\]

Integrating out \(V\),
\[
f_U(u)
=
\int_0^\infty
f_1(v)
f_2\left(\frac vu\right)
\frac{v}{u^2}\,dv,
\qquad u>0.
\]

Therefore,
\[
f_U(u)
=
\begin{cases}
\displaystyle
\int_0^\infty
\frac{v}{u^2}
f_1(v)
f_2\left(\frac vu\right)\,dv,
&
u>0,\\[4mm]
0,
&
u\leq0.
\end{cases}
\]

\end{solution}


%%%%%%%%%%%%%%%%%%%%%%%%%%%%%%%%%%%%%%%%%%%%%%%%%%%%%%%%%%%%%%
%%%%%%%%%%%%%%%%%%%%%%  PROBLEM 18  %%%%%%%%%%%%%%%%%%%%%%%%%%
%%%%%%%%%%%%%%%%%%%%%%%%%%%%%%%%%%%%%%%%%%%%%%%%%%%%%%%%%%%%%%

\newpage
\begin{problem}

Let \(X\) and \(Y\) be independent random variables with common probability
density function
\[
f(x)
=
\begin{cases}
\dfrac1{x^2},
&
x\geq1,\\[2mm]
0,
&
\text{elsewhere}.
\end{cases}
\]

Calculate the joint probability density function of
\[
U=\frac XY
\qquad\text{and}\qquad
V=XY.
\]

\end{problem}

\begin{solution}

The transformation is
\[
u=\frac xy,
\qquad
v=xy.
\]

Since \(x>0\) and \(y>0\),
\[
uv=x^2,
\]
and
\[
\frac vu=y^2.
\]

Therefore,
\[
x=\sqrt{uv},
\qquad
y=\sqrt{\frac vu}.
\]

The Jacobian matrix of the inverse transformation is
\[
\frac{\partial(x,y)}{\partial(u,v)}.
\]

Its absolute determinant is
\[
\left|
\det
\frac{\partial(x,y)}{\partial(u,v)}
\right|
=
\frac1{2u}.
\]

Since \(X\) and \(Y\) are independent,
\[
f_{X,Y}(x,y)
=
\frac1{x^2y^2},
\qquad x\geq1,\quad y\geq1.
\]

But
\[
x^2y^2=(xy)^2=v^2.
\]

Therefore,
\[
f_{X,Y}(x,y)=\frac1{v^2}.
\]

It remains to determine the support.

The condition \(x\geq1\) gives
\[
\sqrt{uv}\geq1,
\]
so
\[
uv\geq1
\]
and therefore
\[
v\geq\frac1u.
\]

The condition \(y\geq1\) gives
\[
\sqrt{\frac vu}\geq1,
\]
so
\[
v\geq u.
\]

Thus,
\[
u>0,
\qquad
v\geq\max\left\{u,\frac1u\right\}.
\]

Hence,
\[
f_{U,V}(u,v)
=
\begin{cases}
\dfrac{1}{2uv^2},
&
u>0,\quad
v\geq\max\left\{u,\dfrac1u\right\},\\[3mm]
0,
&
\text{otherwise}.
\end{cases}
\]

\end{solution}


%%%%%%%%%%%%%%%%%%%%%%%%%%%%%%%%%%%%%%%%%%%%%%%%%%%%%%%%%%%%%%
%%%%%%%%%%%%%%%%%%%%%%  PROBLEM 19  %%%%%%%%%%%%%%%%%%%%%%%%%%
%%%%%%%%%%%%%%%%%%%%%%%%%%%%%%%%%%%%%%%%%%%%%%%%%%%%%%%%%%%%%%

\newpage
\begin{problem}
Let \(X\) and \(Y\) be independent random variables with common probability
density function
\[
f(x)
=
\begin{cases}
e^{-x},
&
x>0,\\
0,
&
\text{elsewhere}.
\end{cases}
\]

Find the joint probability density function of
\[
U=X+Y
\qquad\text{and}\qquad
V=e^X.
\]

\end{problem}

\begin{solution}

The transformation is
\[
u=x+y,
\qquad
v=e^x.
\]

From $v=e^x,$ we obtain  $x=\ln v.$ Therefore, $y=u-\ln v.$ Since \(x>0\), $\ln v>0,$  so $ v>1.$ Since \(y>0\), $u-\ln v>0,$
so $u>\ln v.$
 
The Jacobian matrix of the inverse transformation is
\[
\frac{\partial(x,y)}{\partial(u,v)}
=
\begin{pmatrix}
0 & \dfrac1v\\[2mm]
1 & -\dfrac1v
\end{pmatrix}.
\]

Therefore,
\[
\left|
\det
\frac{\partial(x,y)}{\partial(u,v)}
\right|
=
\frac1v.
\]

Because \(X\) and \(Y\) are independent,
\[
f_{X,Y}(x,y)
=
e^{-x}e^{-y}
=
e^{-(x+y)}.
\]

Since
\[
x+y=u,
\]
we have
\[
f_{X,Y}(x,y)=e^{-u}.
\]

Hence,
\[
f_{U,V}(u,v)
=
\begin{cases}
\dfrac{e^{-u}}{v},
&
v>1,\quad u>\ln v,\\[2mm]
0,
&
\text{otherwise}.
\end{cases}
\]

\end{solution}


%%%%%%%%%%%%%%%%%%%%%%%%%%%%%%%%%%%%%%%%%%%%%%%%%%%%%%%%%%%%%%
%%%%%%%%%%%%%%%%%%%%%%  PROBLEM 20  %%%%%%%%%%%%%%%%%%%%%%%%%%
%%%%%%%%%%%%%%%%%%%%%%%%%%%%%%%%%%%%%%%%%%%%%%%%%%%%%%%%%%%%%%

\newpage
\begin{problem}

Suppose that \(X\) and \(Y\) are independent standard Normal random variables.
Prove that
\[
X+Y
\qquad\text{and}\qquad
X-Y
\]
are independent random variables.

\end{problem}

\begin{solution}

Define
\[
U=X+Y,
\qquad
V=X-Y.
\]

Solving for \(X\) and \(Y\),
\[
X=\frac{U+V}{2},
\qquad
Y=\frac{U-V}{2}.
\]

The Jacobian matrix of the inverse transformation is
\[
\frac{\partial(x,y)}{\partial(u,v)}
=
\begin{pmatrix}
\dfrac12 & \dfrac12\\[2mm]
\dfrac12 & -\dfrac12
\end{pmatrix}.
\]

Therefore,
\[
\left|
\det
\frac{\partial(x,y)}{\partial(u,v)}
\right|
=
\frac12.
\]

Since \(X\) and \(Y\) are independent standard Normal random variables,
\[
f_{X,Y}(x,y)
=
\frac1{2\pi}
\exp\left(
-\frac{x^2+y^2}{2}
\right).
\]

Now
\[
\begin{aligned}
x^2+y^2
&=
\left(\frac{u+v}{2}\right)^2
+
\left(\frac{u-v}{2}\right)^2\\
&=
\frac{u^2+v^2}{2}.
\end{aligned}
\]

Therefore,
\[
\begin{aligned}
f_{U,V}(u,v)
&=
\frac1{2\pi}
\exp\left(
-\frac{u^2+v^2}{4}
\right)
\frac12\\
&=
\frac1{4\pi}
e^{-u^2/4}
e^{-v^2/4}.
\end{aligned}
\]

But
\[
\frac1{\sqrt{4\pi}}e^{-u^2/4}
\]
is the density of a \(N(0,2)\) random variable. Hence,
\[
f_U(u)
=
\frac1{\sqrt{4\pi}}e^{-u^2/4}
\]
and
\[
f_V(v)
=
\frac1{\sqrt{4\pi}}e^{-v^2/4}.
\]

Consequently, $f_{U,V}(u,v)
=
f_U(u)f_V(v)$ and thus the independence.
 


\end{solution}


%%%%%%%%%%%%%%%%%%%%%%%%%%%%%%%%%%%%%%%%%%%%%%%%%%%%%%%%%%%%%%
%%%%%%%%%%%%%%%%%%%%%%  PROBLEM 21  %%%%%%%%%%%%%%%%%%%%%%%%%%
%%%%%%%%%%%%%%%%%%%%%%%%%%%%%%%%%%%%%%%%%%%%%%%%%%%%%%%%%%%%%%

\newpage
\begin{problem}

Let \(X\) and \(Y\) be independent strictly positive exponential random
variables, each with parameter \(\lambda\). Are the random variables
\[
X+Y
\qquad\text{and}\qquad
\frac XY
\]
independent?

\end{problem}

\begin{solution}

Define
\[
U=X+Y,
\qquad
V=\frac XY.
\]

Since
\[
v=\frac xy,
\]
we have
\[
x=vy.
\]

Substituting into
\[
u=x+y
\]
gives
\[
u=(v+1)y.
\]

Therefore,
\[
y=\frac{u}{1+v}
\]
and
\[
x=\frac{uv}{1+v}.
\]

Since \(X,Y>0\),
\[
u>0,
\qquad
v>0.
\]

The Jacobian matrix is
\[
\frac{\partial(x,y)}{\partial(u,v)}
=
\begin{pmatrix}
\dfrac{v}{1+v}
&
\dfrac{u}{(1+v)^2}
\\[3mm]
\dfrac1{1+v}
&
-\dfrac{u}{(1+v)^2}
\end{pmatrix}.
\]

Therefore,
\[
\left|
\det
\frac{\partial(x,y)}{\partial(u,v)}
\right|
=
\frac{u}{(1+v)^2}.
\]

Because \(X\) and \(Y\) are independent exponential random variables,
\[
f_{X,Y}(x,y)
=
\lambda^2e^{-\lambda(x+y)},
\qquad x,y>0.
\]

Since
\[
x+y=u,
\]
the transformed joint density is
\[
\begin{aligned}
f_{U,V}(u,v)
&=
\lambda^2e^{-\lambda u}
\frac{u}{(1+v)^2}\\
&=
\left(
\lambda^2u e^{-\lambda u}
\right)
\left(
\frac1{(1+v)^2}
\right),
\end{aligned}
\]
for
\[
u>0,
\qquad
v>0.
\]

The first factor is a density on \(u>0\), since
\[
\int_0^\infty
\lambda^2u e^{-\lambda u}\,du
=
1.
\]

The second factor is also a density on \(v>0\), since
\[
\int_0^\infty
\frac1{(1+v)^2}\,dv
=
1.
\]

Thus,
\[
f_U(u)
=
\lambda^2u e^{-\lambda u},
\qquad u>0,
\]
and
\[
f_V(v)
=
\frac1{(1+v)^2},
\qquad v>0.
\]

Therefore,
\[
f_{U,V}(u,v)
=
f_U(u)f_V(v).
\]

Hence,
\[
X+Y
\qquad\text{and}\qquad
\frac XY
\]
are independent.

\end{solution}




 






%%%%%%%%%%%%%%%%%%%%%%%%%%%%%%%%%%%%%%%%%%%%%%%%%%%%%%%%%%%%%%
%%%%%%%%%%%%%%%%%%%%%%  PROBLEM 22  %%%%%%%%%%%%%%%%%%%%%%%%%%
%%%%%%%%%%%%%%%%%%%%%%%%%%%%%%%%%%%%%%%%%%%%%%%%%%%%%%%%%%%%%%

\newpage
\begin{problem}

Let
\[
p(x,y,z)=\frac{xyz}{162},
\qquad
x=4,5,\quad
y=1,2,3,\quad
z=1,2,
\]
be the joint probability mass function of random variables \(X,Y,Z\).

\begin{enumerate}[label=(\alph*)]
    \item Calculate the joint marginal probability mass functions of
    \((X,Y)\), \((Y,Z)\), and \((X,Z)\).
    \item Find \(E(YZ)\).
\end{enumerate}

\end{problem}

\begin{solution}

\begin{enumerate}[label=(\alph*)]

\item To obtain the joint marginal PMF of \(X\) and \(Y\), sum over all
possible values of \(Z\):
\[
\begin{aligned}
p_{X,Y}(x,y)
&=
\sum_{z=1}^{2}p(x,y,z)\\
&=
\sum_{z=1}^{2}\frac{xyz}{162}\\
&=
\frac{xy}{162}(1+2)\\
&=
\frac{xy}{54}.
\end{aligned}
\]

Thus,
\[
p_{X,Y}(x,y)
=
\begin{cases}
\dfrac{xy}{54},
&
x=4,5,\quad y=1,2,3,\\[2mm]
0,
&
\text{otherwise}.
\end{cases}
\]

Similarly,
\[
\begin{aligned}
p_{Y,Z}(y,z)
&=
\sum_{x=4}^{5}p(x,y,z)\\
&=
\frac{yz}{162}(4+5)\\
&=
\frac{yz}{18}.
\end{aligned}
\]

Therefore,
\[
p_{Y,Z}(y,z)
=
\begin{cases}
\dfrac{yz}{18},
&
y=1,2,3,\quad z=1,2,\\[2mm]
0,
&
\text{otherwise}.
\end{cases}
\]

Finally,
\[
\begin{aligned}
p_{X,Z}(x,z)
&=
\sum_{y=1}^{3}p(x,y,z)\\
&=
\frac{xz}{162}(1+2+3)\\
&=
\frac{xz}{27}.
\end{aligned}
\]

Hence,
\[
p_{X,Z}(x,z)
=
\begin{cases}
\dfrac{xz}{27},
&
x=4,5,\quad z=1,2,\\[2mm]
0,
&
\text{otherwise}.
\end{cases}
\]

\item Using the joint marginal PMF of \(Y\) and \(Z\),
\[
E(YZ)
=
\sum_{y=1}^{3}\sum_{z=1}^{2}
yz\,p_{Y,Z}(y,z).
\]

Thus,
\[
\begin{aligned}
E(YZ)
&=
\sum_{y=1}^{3}\sum_{z=1}^{2}
yz\frac{yz}{18}\\
&=
\frac1{18}
\left(\sum_{y=1}^{3}y^2\right)
\left(\sum_{z=1}^{2}z^2\right)\\
&=
\frac1{18}(1+4+9)(1+4)\\
&=
\frac{70}{18}\\
&=
\frac{35}{9}.
\end{aligned}
\]

Therefore,
\[
E(YZ)=\frac{35}{9}.
\]

\end{enumerate}

\end{solution}


%%%%%%%%%%%%%%%%%%%%%%%%%%%%%%%%%%%%%%%%%%%%%%%%%%%%%%%%%%%%%%
%%%%%%%%%%%%%%%%%%%%%%  PROBLEM 23  %%%%%%%%%%%%%%%%%%%%%%%%%%
%%%%%%%%%%%%%%%%%%%%%%%%%%%%%%%%%%%%%%%%%%%%%%%%%%%%%%%%%%%%%%

\newpage
\begin{problem}

Let the joint probability density function of \(X,Y,Z\) be
\[
f(x,y,z)
=
\begin{cases}
6e^{-x-y-z},
&
0<x<y<z<\infty,\\
0,
&
\text{elsewhere}.
\end{cases}
\]

\begin{enumerate}[label=(\alph*)]
    \item Find the joint marginal probability density functions of
    \((X,Y)\), \((X,Z)\), and \((Y,Z)\).
    \item Find \(E(X)\).
\end{enumerate}

\end{problem}

\begin{solution}

The support is
\[
0<x<y<z<\infty.
\]

\begin{enumerate}[label=(\alph*)]

\item For fixed \(x\) and \(y\) satisfying \(0<x<y\),
\[
z>y.
\]

Therefore,
\[
\begin{aligned}
f_{X,Y}(x,y)
&=
\int_y^\infty
6e^{-x-y-z}\,dz\\
&=
6e^{-x-y}
\int_y^\infty e^{-z}\,dz\\
&=
6e^{-x-2y}.
\end{aligned}
\]

Hence,
\[
f_{X,Y}(x,y)
=
\begin{cases}
6e^{-x-2y}, & 0<x<y,\\
0, & \text{otherwise}.
\end{cases}
\]

For fixed \(x\) and \(z\) satisfying \(0<x<z\),
\[
x<y<z.
\]

Thus,
\[
\begin{aligned}
f_{X,Z}(x,z)
&=
\int_x^z
6e^{-x-y-z}\,dy\\
&=
6e^{-x-z}
\left(e^{-x}-e^{-z}\right)\\
&=
6\left(
e^{-2x-z}-e^{-x-2z}
\right).
\end{aligned}
\]

Therefore,
\[
f_{X,Z}(x,z)
=
\begin{cases}
6\left(e^{-2x-z}-e^{-x-2z}\right),
&
0<x<z,\\
0,
&
\text{otherwise}.
\end{cases}
\]

For fixed \(y\) and \(z\) satisfying \(0<y<z\),
\[
0<x<y.
\]

Hence,
\[
\begin{aligned}
f_{Y,Z}(y,z)
&=
\int_0^y
6e^{-x-y-z}\,dx\\
&=
6e^{-y-z}(1-e^{-y}).
\end{aligned}
\]

Thus,
\[
f_{Y,Z}(y,z)
=
\begin{cases}
6e^{-y-z}(1-e^{-y}),
&
0<y<z,\\
0,
&
\text{otherwise}.
\end{cases}
\]

\item We first find the marginal density of \(X\):
\[
\begin{aligned}
f_X(x)
&=
\int_x^\infty
\int_y^\infty
6e^{-x-y-z}\,dz\,dy\\
&=
6e^{-x}
\int_x^\infty
e^{-y}
\left(
\int_y^\infty e^{-z}\,dz
\right)dy\\
&=
6e^{-x}
\int_x^\infty e^{-2y}\,dy\\
&=
6e^{-x}
\left(\frac12e^{-2x}\right)\\
&=
3e^{-3x},
\qquad x>0.
\end{aligned}
\]

Therefore,
\[
\begin{aligned}
E(X)
&=
\int_0^\infty
x\,3e^{-3x}\,dx\\
&=
\frac13.
\end{aligned}
\]

Thus,
\[
E(X)=\frac13.
\]

\end{enumerate}

\end{solution}


%%%%%%%%%%%%%%%%%%%%%%%%%%%%%%%%%%%%%%%%%%%%%%%%%%%%%%%%%%%%%%
%%%%%%%%%%%%%%%%%%%%%%  PROBLEM 24  %%%%%%%%%%%%%%%%%%%%%%%%%%
%%%%%%%%%%%%%%%%%%%%%%%%%%%%%%%%%%%%%%%%%%%%%%%%%%%%%%%%%%%%%%

\newpage
\begin{problem}

Let \(X,Y,Z\) be jointly continuous with joint probability density function
\[
f(x,y,z)
=
\begin{cases}
x^2e^{-x(1+y+z)},
&
x,y,z>0,\\
0,
&
\text{otherwise}.
\end{cases}
\]

Are \(X,Y,Z\) mutually independent? Are they pairwise independent?

\end{problem}

\begin{solution}

We first calculate the marginal densities.

For \(x>0\),
\[
\begin{aligned}
f_X(x)
&=
\int_0^\infty
\int_0^\infty
x^2e^{-x(1+y+z)}
\,dy\,dz\\
&=
x^2e^{-x}
\left(
\int_0^\infty e^{-xy}\,dy
\right)
\left(
\int_0^\infty e^{-xz}\,dz
\right)\\
&=
x^2e^{-x}
\left(\frac1x\right)
\left(\frac1x\right)\\
&=
e^{-x}.
\end{aligned}
\]

For \(y>0\),
\[
\begin{aligned}
f_Y(y)
&=
\int_0^\infty
\int_0^\infty
x^2e^{-x(1+y+z)}
\,dz\,dx\\
&=
\int_0^\infty
xe^{-x(1+y)}\,dx\\
&=
\frac{1}{(1+y)^2}.
\end{aligned}
\]

By symmetry,
\[
f_Z(z)=\frac{1}{(1+z)^2},
\qquad z>0.
\]

If \(X,Y,Z\) were mutually independent, then
\[
f(x,y,z)
=
f_X(x)f_Y(y)f_Z(z).
\]

However,
\[
f_X(x)f_Y(y)f_Z(z)
=
\frac{e^{-x}}
{(1+y)^2(1+z)^2},
\]
which is not equal to
\[
x^2e^{-x(1+y+z)}.
\]

Therefore, the three random variables are not mutually independent.

We now examine pairwise independence.

For \(x,y>0\),
\[
\begin{aligned}
f_{X,Y}(x,y)
&=
\int_0^\infty
x^2e^{-x(1+y+z)}\,dz\\
&=
xe^{-x(1+y)}.
\end{aligned}
\]

But
\[
f_X(x)f_Y(y)
=
\frac{e^{-x}}{(1+y)^2},
\]
which is not equal to \(f_{X,Y}(x,y)\). Hence \(X\) and \(Y\) are not
independent.

Similarly, \(X\) and \(Z\) are not independent.

Finally,
\[
\begin{aligned}
f_{Y,Z}(y,z)
&=
\int_0^\infty
x^2e^{-x(1+y+z)}\,dx\\
&=
\frac{2}{(1+y+z)^3},
\end{aligned}
\]
whereas
\[
f_Y(y)f_Z(z)
=
\frac{1}{(1+y)^2(1+z)^2}.
\]

Thus \(Y\) and \(Z\) are not independent either.

Therefore, \(X,Y,Z\) are neither mutually independent nor pairwise
independent.

\end{solution}


%%%%%%%%%%%%%%%%%%%%%%%%%%%%%%%%%%%%%%%%%%%%%%%%%%%%%%%%%%%%%%
%%%%%%%%%%%%%%%%%%%%%%  PROBLEM 25  %%%%%%%%%%%%%%%%%%%%%%%%%%
%%%%%%%%%%%%%%%%%%%%%%%%%%%%%%%%%%%%%%%%%%%%%%%%%%%%%%%%%%%%%%

\newpage
\begin{problem}

Let the joint cumulative distribution function of \(X,Y,Z\) be
\[
F(x,y,z)
=
\left(1-e^{-\lambda_1x}\right)
\left(1-e^{-\lambda_2y}\right)
\left(1-e^{-\lambda_3z}\right),
\qquad x,y,z>0,
\]
where
\[
\lambda_1,\lambda_2,\lambda_3>0.
\]

\begin{enumerate}[label=(\alph*)]
    \item Are \(X,Y,Z\) independent?
    \item Find the joint probability density function of \(X,Y,Z\).
    \item Find \(P(X<Y<Z)\).
\end{enumerate}

\end{problem}

\begin{solution}

\begin{enumerate}[label=(\alph*)]

\item The joint CDF factors as
\[
F(x,y,z)
=
F_X(x)F_Y(y)F_Z(z),
\]
where
\[
F_X(x)=1-e^{-\lambda_1x},
\]
\[
F_Y(y)=1-e^{-\lambda_2y},
\]
and
\[
F_Z(z)=1-e^{-\lambda_3z},
\]
for positive arguments.

Therefore,
\[
X\sim\operatorname{Exp}(\lambda_1),
\qquad
Y\sim\operatorname{Exp}(\lambda_2),
\qquad
Z\sim\operatorname{Exp}(\lambda_3),
\]
and \(X,Y,Z\) are mutually independent.

\item Differentiating the joint CDF,
\[
f(x,y,z)
=
\frac{\partial^3}
{\partial x\,\partial y\,\partial z}
F(x,y,z).
\]

Therefore,
\[
f(x,y,z)
=
\lambda_1\lambda_2\lambda_3
e^{-\lambda_1x-\lambda_2y-\lambda_3z},
\qquad
x,y,z>0.
\]

Thus,
\[
f(x,y,z)
=
\begin{cases}
\lambda_1\lambda_2\lambda_3
e^{-\lambda_1x-\lambda_2y-\lambda_3z},
&
x,y,z>0,\\
0,
&
\text{otherwise}.
\end{cases}
\]

\item We calculate
\[
P(X<Y<Z)
=
\int_0^\infty
\int_x^\infty
\int_y^\infty
f(x,y,z)\,dz\,dy\,dx.
\]

Hence,
\[
\begin{aligned}
P(X<Y<Z)
&=
\int_0^\infty
\int_x^\infty
\int_y^\infty
\lambda_1\lambda_2\lambda_3
e^{-\lambda_1x-\lambda_2y-\lambda_3z}
\,dz\,dy\,dx\\
&=
\int_0^\infty
\int_x^\infty
\lambda_1\lambda_2
e^{-\lambda_1x-(\lambda_2+\lambda_3)y}
\,dy\,dx\\
&=
\frac{\lambda_1\lambda_2}
{\lambda_2+\lambda_3}
\int_0^\infty
e^{-(\lambda_1+\lambda_2+\lambda_3)x}
\,dx\\
&=
\frac{\lambda_1\lambda_2}
{(\lambda_2+\lambda_3)
(\lambda_1+\lambda_2+\lambda_3)}.
\end{aligned}
\]

Therefore,
\[
P(X<Y<Z)
=
\frac{\lambda_1\lambda_2}
{(\lambda_1+\lambda_2+\lambda_3)(\lambda_2+\lambda_3)}.
\]

\end{enumerate}

\end{solution}


%%%%%%%%%%%%%%%%%%%%%%%%%%%%%%%%%%%%%%%%%%%%%%%%%%%%%%%%%%%%%%
%%%%%%%%%%%%%%%%%%%%%%  PROBLEM 26  %%%%%%%%%%%%%%%%%%%%%%%%%%
%%%%%%%%%%%%%%%%%%%%%%%%%%%%%%%%%%%%%%%%%%%%%%%%%%%%%%%%%%%%%%

\newpage
\begin{problem}

Suppose that \(h\) is the probability density function of a continuous random
variable. Let the joint probability density function of \(X,Y,Z\) be
\[
f(x,y,z)=h(x)h(y)h(z),
\qquad
x,y,z\in\mathbb R.
\]

Prove that
\[
P(X<Y<Z)=\frac16.
\]

\end{problem}

\begin{solution}

Since
\[
f(x,y,z)=h(x)h(y)h(z),
\]
the random variables \(X,Y,Z\) are independent and identically distributed.

Because their common distribution is continuous, ties occur with probability
zero:
\[
P(X=Y)=P(X=Z)=P(Y=Z)=0.
\]

Therefore, with probability \(1\), the three values occur in exactly one of
the six possible strict orderings:
\[
X<Y<Z,
\]
\[
X<Z<Y,
\]
\[
Y<X<Z,
\]
\[
Y<Z<X,
\]
\[
Z<X<Y,
\]
or
\[
Z<Y<X.
\]

Since \(X,Y,Z\) are identically distributed, symmetry implies that all six
orderings have the same probability.

Their probabilities sum to \(1\), so each must have probability
\[
\frac16.
\]

Hence,
\[
P(X<Y<Z)=\frac16.
\]

\end{solution}


%%%%%%%%%%%%%%%%%%%%%%%%%%%%%%%%%%%%%%%%%%%%%%%%%%%%%%%%%%%%%%
%%%%%%%%%%%%%%%%%%%%%%  PROBLEM 27  %%%%%%%%%%%%%%%%%%%%%%%%%%
%%%%%%%%%%%%%%%%%%%%%%%%%%%%%%%%%%%%%%%%%%%%%%%%%%%%%%%%%%%%%%

\newpage
\begin{problem}

Let \(X\) and \(Y\) be random variables with finite second moments, and let
\(a,b\) be constants. Prove that
\[
\operatorname{Var}(aX+bY)
=
a^2\operatorname{Var}(X)
+
b^2\operatorname{Var}(Y)
+
2ab\operatorname{Cov}(X,Y).
\]

Deduce that if \(X\) and \(Y\) are independent, then
\[
\operatorname{Var}(aX+bY)
=
a^2\operatorname{Var}(X)
+
b^2\operatorname{Var}(Y).
\]

\end{problem}

\begin{solution}

By definition,
\[
\operatorname{Var}(aX+bY)
=
E\left[
(aX+bY)-E(aX+bY)
\right]^2.
\]

By linearity of expectation,
\[
E(aX+bY)=aE(X)+bE(Y).
\]

Therefore,
\[
\begin{aligned}
\operatorname{Var}(aX+bY)
&=
E\left[
a(X-E(X))
+
b(Y-E(Y))
\right]^2\\
&=
E\Big[
a^2(X-E(X))^2
+
b^2(Y-E(Y))^2\\
&\hspace{2.7cm}
+
2ab(X-E(X))(Y-E(Y))
\Big].
\end{aligned}
\]

Using linearity of expectation,
\[
\begin{aligned}
\operatorname{Var}(aX+bY)
&=
a^2E[(X-E(X))^2]
+
b^2E[(Y-E(Y))^2]\\
&\qquad
+
2abE[(X-E(X))(Y-E(Y))].
\end{aligned}
\]

Thus,
\[
\operatorname{Var}(aX+bY)
=
a^2\operatorname{Var}(X)
+
b^2\operatorname{Var}(Y)
+
2ab\operatorname{Cov}(X,Y).
\]

If \(X\) and \(Y\) are independent, then
\[
\operatorname{Cov}(X,Y)=0.
\]

Hence,
\[
\operatorname{Var}(aX+bY)
=
a^2\operatorname{Var}(X)
+
b^2\operatorname{Var}(Y).
\]

\end{solution}


%%%%%%%%%%%%%%%%%%%%%%%%%%%%%%%%%%%%%%%%%%%%%%%%%%%%%%%%%%%%%%
%%%%%%%%%%%%%%%%%%%%%%  PROBLEM 28  %%%%%%%%%%%%%%%%%%%%%%%%%%
%%%%%%%%%%%%%%%%%%%%%%%%%%%%%%%%%%%%%%%%%%%%%%%%%%%%%%%%%%%%%%

\newpage
\begin{problem}

Let \(X_1,\ldots,X_n\) be random variables with finite second moments, and let
\(a_1,\ldots,a_n\) be constants. Prove that
\[
\operatorname{Var}
\left(
\sum_{i=1}^{n}a_iX_i
\right)
=
\sum_{i=1}^{n}
a_i^2\operatorname{Var}(X_i)
+
2\sum_{i<j}
a_ia_j\operatorname{Cov}(X_i,X_j).
\]

In particular, show that
\[
\operatorname{Var}
\left(
\sum_{i=1}^{n}X_i
\right)
=
\sum_{i=1}^{n}
\operatorname{Var}(X_i)
+
2\sum_{i<j}
\operatorname{Cov}(X_i,X_j).
\]

\end{problem}

\begin{solution}

Let
\[
S=\sum_{i=1}^{n}a_iX_i.
\]

Then
\[
E(S)
=
\sum_{i=1}^{n}a_iE(X_i).
\]

Therefore,
\[
S-E(S)
=
\sum_{i=1}^{n}
a_i\bigl(X_i-E(X_i)\bigr).
\]

Hence,
\[
\operatorname{Var}(S)
=
E\left[
\left(
\sum_{i=1}^{n}
a_i(X_i-E(X_i))
\right)^2
\right].
\]

Expanding the square,
\[
\begin{aligned}
\operatorname{Var}(S)
&=
E\left[
\sum_{i=1}^{n}
a_i^2(X_i-E(X_i))^2
\right.\\
&\qquad\left.
+
2\sum_{i<j}
a_ia_j
(X_i-E(X_i))(X_j-E(X_j))
\right].
\end{aligned}
\]

By linearity of expectation,
\[
\operatorname{Var}(S)
=
\sum_{i=1}^{n}
a_i^2\operatorname{Var}(X_i)
+
2\sum_{i<j}
a_ia_j\operatorname{Cov}(X_i,X_j).
\]

Taking
\[
a_i=1,
\qquad i=1,\ldots,n,
\]
gives
\[
\operatorname{Var}
\left(
\sum_{i=1}^{n}X_i
\right)
=
\sum_{i=1}^{n}
\operatorname{Var}(X_i)
+
2\sum_{i<j}
\operatorname{Cov}(X_i,X_j).
\]

\end{solution}


%%%%%%%%%%%%%%%%%%%%%%%%%%%%%%%%%%%%%%%%%%%%%%%%%%%%%%%%%%%%%%
%%%%%%%%%%%%%%%%%%%%%%  PROBLEM 29  %%%%%%%%%%%%%%%%%%%%%%%%%%
%%%%%%%%%%%%%%%%%%%%%%%%%%%%%%%%%%%%%%%%%%%%%%%%%%%%%%%%%%%%%%

\newpage
\begin{problem}

Let \(X_1,\ldots,X_n\) be a random sample from a distribution with mean
\(\mu\) and variance \(\sigma^2\). Define
\[
\overline X
=
\frac1n\sum_{i=1}^{n}X_i.
\]

Prove that
\[
E(\overline X)=\mu
\]
and
\[
\operatorname{Var}(\overline X)
=
\frac{\sigma^2}{n}.
\]

\end{problem}

\begin{solution}

Since \(X_1,\ldots,X_n\) form a random sample,
\[
E(X_i)=\mu
\]
and
\[
\operatorname{Var}(X_i)=\sigma^2
\]
for every \(i\), and the random variables are independent.

By linearity of expectation,
\[
\begin{aligned}
E(\overline X)
&=
E\left(
\frac1n\sum_{i=1}^{n}X_i
\right)\\
&=
\frac1n
\sum_{i=1}^{n}E(X_i)\\
&=
\frac1n(n\mu)\\
&=
\mu.
\end{aligned}
\]

For the variance, independence implies
\[
\operatorname{Cov}(X_i,X_j)=0
\qquad
(i\neq j).
\]

Hence,
\[
\begin{aligned}
\operatorname{Var}(\overline X)
&=
\operatorname{Var}
\left(
\frac1n\sum_{i=1}^{n}X_i
\right)\\
&=
\frac1{n^2}
\sum_{i=1}^{n}
\operatorname{Var}(X_i)\\
&=
\frac1{n^2}(n\sigma^2)\\
&=
\frac{\sigma^2}{n}.
\end{aligned}
\]

Thus,
\[
E(\overline X)=\mu,
\qquad
\operatorname{Var}(\overline X)=\frac{\sigma^2}{n}.
\]

\end{solution}


%%%%%%%%%%%%%%%%%%%%%%%%%%%%%%%%%%%%%%%%%%%%%%%%%%%%%%%%%%%%%%
%%%%%%%%%%%%%%%%%%%%%%  PROBLEM 30  %%%%%%%%%%%%%%%%%%%%%%%%%%
%%%%%%%%%%%%%%%%%%%%%%%%%%%%%%%%%%%%%%%%%%%%%%%%%%%%%%%%%%%%%%

\newpage
\begin{problem}

Let \(X\) be uniformly distributed on \((-1,1)\), and define
\[
Y=X^2.
\]

Show that \(X\) and \(Y\) are dependent but uncorrelated.

\end{problem}

\begin{solution}

Since
\[
Y=X^2,
\]
the value of \(Y\) is completely determined by \(X\). Therefore, \(X\) and
\(Y\) are not independent.

We now calculate their covariance.

By symmetry of the uniform distribution on \((-1,1)\),
\[
E(X)=0.
\]

Also,
\[
E(X^3)=0,
\]
because \(x^3\) is an odd function and the distribution of \(X\) is symmetric
about zero.

Since
\[
Y=X^2,
\]
we have
\[
XY=X^3.
\]

Thus,
\[
\begin{aligned}
\operatorname{Cov}(X,Y)
&=
E(XY)-E(X)E(Y)\\
&=
E(X^3)-E(X)E(X^2)\\
&=
0-0\cdot E(X^2)\\
&=
0.
\end{aligned}
\]

Therefore, \(X\) and \(Y\) are dependent but uncorrelated.

\end{solution}


%%%%%%%%%%%%%%%%%%%%%%%%%%%%%%%%%%%%%%%%%%%%%%%%%%%%%%%%%%%%%%
%%%%%%%%%%%%%%%%%%%%%%  PROBLEM 31  %%%%%%%%%%%%%%%%%%%%%%%%%%
%%%%%%%%%%%%%%%%%%%%%%%%%%%%%%%%%%%%%%%%%%%%%%%%%%%%%%%%%%%%%%

\newpage
\begin{problem}

Let the joint probability mass function of random variables \(X\) and \(Y\)
be
\[
p(x,y)
=
\begin{cases}
\dfrac1{70}x(x+y),
&
x=1,2,3,\quad y=3,4,\\[2mm]
0,
&
\text{elsewhere}.
\end{cases}
\]

Find
\[
\operatorname{Cov}(X,Y).
\]

\end{problem}

\begin{solution}

Recall that
\[
\operatorname{Cov}(X,Y)
=
E(XY)-E(X)E(Y).
\]

First calculate \(E(X)\):
\[
\begin{aligned}
E(X)
&=
\sum_{x=1}^{3}
\sum_{y=3}^{4}
x\,p(x,y)\\
&=
\frac1{70}
\sum_{x=1}^{3}
\sum_{y=3}^{4}
x^2(x+y).
\end{aligned}
\]

Evaluating the six terms,
\[
E(X)
=
\frac{156}{70}
=
\frac{78}{35}.
\]

Next,
\[
\begin{aligned}
E(Y)
&=
\sum_{x=1}^{3}
\sum_{y=3}^{4}
y\,p(x,y)\\
&=
\frac1{70}
\sum_{x=1}^{3}
\sum_{y=3}^{4}
xy(x+y)\\
&=
\frac{249}{70}.
\end{aligned}
\]

Finally,
\[
\begin{aligned}
E(XY)
&=
\sum_{x=1}^{3}
\sum_{y=3}^{4}
xy\,p(x,y)\\
&=
\frac1{70}
\sum_{x=1}^{3}
\sum_{y=3}^{4}
x^2y(x+y)\\
&=
\frac{558}{70}\\
&=
\frac{279}{35}.
\end{aligned}
\]

Therefore,
\[
\begin{aligned}
\operatorname{Cov}(X,Y)
&=
\frac{279}{35}
-
\left(\frac{78}{35}\right)
\left(\frac{249}{70}\right)\\
&=
\frac{279}{35}
-
\frac{9711}{1225}\\
&=
\frac{54}{1225}.
\end{aligned}
\]

 

\end{solution}


%%%%%%%%%%%%%%%%%%%%%%%%%%%%%%%%%%%%%%%%%%%%%%%%%%%%%%%%%%%%%%
%%%%%%%%%%%%%%%%%%%%%%  PROBLEM 32  %%%%%%%%%%%%%%%%%%%%%%%%%%
%%%%%%%%%%%%%%%%%%%%%%%%%%%%%%%%%%%%%%%%%%%%%%%%%%%%%%%%%%%%%%

\newpage
\begin{problem}

Let \(X,Y,Z\) be random variables with finite second moments.

\begin{enumerate}[label=(\alph*)]
    \item Prove that
    \[
    \operatorname{Cov}(X+Y,Z)
    =
    \operatorname{Cov}(X,Z)
    +
    \operatorname{Cov}(Y,Z).
    \]

    \item Prove that
    \[
    \operatorname{Cov}(X,Y+Z)
    =
    \operatorname{Cov}(X,Y)
    +
    \operatorname{Cov}(X,Z).
    \]

    \item Deduce that
    \[
    \operatorname{Cov}(X+Y,X-Y)
    =
    \operatorname{Var}(X)-\operatorname{Var}(Y).
    \]

    \item Deduce that
    \[
    \operatorname{Var}(X-Y)
    =
    \operatorname{Var}(X)
    +
    \operatorname{Var}(Y)
    -
    2\operatorname{Cov}(X,Y).
    \]
\end{enumerate}

\end{problem}

\begin{solution}

Recall that
\[
\operatorname{Cov}(A,B)
=
E(AB)-E(A)E(B).
\]

\begin{enumerate}[label=(\alph*)]

\item
\[
\begin{aligned}
\operatorname{Cov}(X+Y,Z)
&=
E[(X+Y)Z]
-
E(X+Y)E(Z)\\
&=
E(XZ)+E(YZ)
-
[E(X)+E(Y)]E(Z)\\
&=
[E(XZ)-E(X)E(Z)]\\
&\qquad+
[E(YZ)-E(Y)E(Z)]\\
&=
\operatorname{Cov}(X,Z)
+
\operatorname{Cov}(Y,Z).
\end{aligned}
\]

\item Similarly,
\[
\begin{aligned}
\operatorname{Cov}(X,Y+Z)
&=
E[X(Y+Z)]
-
E(X)E(Y+Z)\\
&=
\operatorname{Cov}(X,Y)
+
\operatorname{Cov}(X,Z).
\end{aligned}
\]

\item Using bilinearity,
\[
\begin{aligned}
\operatorname{Cov}(X+Y,X-Y)
&=
\operatorname{Cov}(X,X)
-
\operatorname{Cov}(X,Y)\\
&\qquad+
\operatorname{Cov}(Y,X)
-
\operatorname{Cov}(Y,Y).
\end{aligned}
\]

Since
\[
\operatorname{Cov}(X,X)=\operatorname{Var}(X),
\]
\[
\operatorname{Cov}(Y,Y)=\operatorname{Var}(Y),
\]
and
\[
\operatorname{Cov}(X,Y)=\operatorname{Cov}(Y,X),
\]
the middle terms cancel. Therefore,
\[
\operatorname{Cov}(X+Y,X-Y)
=
\operatorname{Var}(X)-\operatorname{Var}(Y).
\]

\item Since
\[
\operatorname{Var}(X-Y)
=
\operatorname{Cov}(X-Y,X-Y),
\]
bilinearity gives
\[
\begin{aligned}
\operatorname{Var}(X-Y)
&=
\operatorname{Cov}(X,X)
-
\operatorname{Cov}(X,Y)\\
&\qquad
-
\operatorname{Cov}(Y,X)
+
\operatorname{Cov}(Y,Y)\\
&=
\operatorname{Var}(X)
+
\operatorname{Var}(Y)
-
2\operatorname{Cov}(X,Y).
\end{aligned}
\]

\end{enumerate}

\end{solution}


%%%%%%%%%%%%%%%%%%%%%%%%%%%%%%%%%%%%%%%%%%%%%%%%%%%%%%%%%%%%%%
%%%%%%%%%%%%%%%%%%%%%%  PROBLEM 33  %%%%%%%%%%%%%%%%%%%%%%%%%%
%%%%%%%%%%%%%%%%%%%%%%%%%%%%%%%%%%%%%%%%%%%%%%%%%%%%%%%%%%%%%%

\newpage
\begin{problem}

Let \(X_1,\ldots,X_n\) and \(Y_1,\ldots,Y_m\) be random variables with finite
second moments, and let \(a_1,\ldots,a_n\) and \(b_1,\ldots,b_m\) be constants.

Prove that
\[
\operatorname{Cov}
\left(
\sum_{i=1}^{n}a_iX_i,
\sum_{j=1}^{m}b_jY_j
\right)
=
\sum_{i=1}^{n}
\sum_{j=1}^{m}
a_ib_j\operatorname{Cov}(X_i,Y_j).
\]

Deduce, as a special case, that for constants \(a,b,c,d\),
\[
\operatorname{Cov}(aX+b,cY+d)
=
ac\,\operatorname{Cov}(X,Y).
\]

\end{problem}

\begin{solution}

Using bilinearity of covariance repeatedly,
\[
\begin{aligned}
\operatorname{Cov}
\left(
\sum_{i=1}^{n}a_iX_i,
\sum_{j=1}^{m}b_jY_j
\right)
&=
\sum_{i=1}^{n}
a_i
\operatorname{Cov}
\left(
X_i,
\sum_{j=1}^{m}b_jY_j
\right)\\
&=
\sum_{i=1}^{n}
a_i
\sum_{j=1}^{m}
b_j\operatorname{Cov}(X_i,Y_j)\\
&=
\sum_{i=1}^{n}
\sum_{j=1}^{m}
a_ib_j\operatorname{Cov}(X_i,Y_j).
\end{aligned}
\]

For the special case,
\[
\operatorname{Cov}(aX+b,cY+d),
\]
constants have zero covariance with every random variable. Hence,
\[
\begin{aligned}
\operatorname{Cov}(aX+b,cY+d)
&=
ac\,\operatorname{Cov}(X,Y).
\end{aligned}
\]

\end{solution}


%%%%%%%%%%%%%%%%%%%%%%%%%%%%%%%%%%%%%%%%%%%%%%%%%%%%%%%%%%%%%%
%%%%%%%%%%%%%%%%%%%%%%  PROBLEM 34  %%%%%%%%%%%%%%%%%%%%%%%%%%
%%%%%%%%%%%%%%%%%%%%%%%%%%%%%%%%%%%%%%%%%%%%%%%%%%%%%%%%%%%%%%

\newpage
\begin{problem}

Let \(A\) and \(B\) be events with
\[
P(A)>0,
\qquad
P(B)>0.
\]

Let \(I_A\) and \(I_B\) denote their indicator random variables:
\[
I_A=
\begin{cases}
1, & \text{if }A\text{ occurs},\\
0, & \text{otherwise},
\end{cases}
\]
and
\[
I_B=
\begin{cases}
1, & \text{if }B\text{ occurs},\\
0, & \text{otherwise}.
\end{cases}
\]

Show that \(I_A\) and \(I_B\) are positively correlated if and only if
\[
P(A\mid B)>P(A),
\]
and equivalently if and only if
\[
P(B\mid A)>P(B).
\]

\end{problem}

\begin{solution}

We have
\[
E(I_A)=P(A),
\qquad
E(I_B)=P(B).
\]

Also,
\[
I_AI_B=I_{A\cap B},
\]
so
\[
E(I_AI_B)=P(A\cap B).
\]

Therefore,
\[
\operatorname{Cov}(I_A,I_B)
=
P(A\cap B)-P(A)P(B).
\]

The indicators are positively correlated if and only if
\[
\operatorname{Cov}(I_A,I_B)>0.
\]

Thus,
\[
P(A\cap B)>P(A)P(B).
\]

Since \(P(B)>0\), dividing by \(P(B)\) gives
\[
\frac{P(A\cap B)}{P(B)}>P(A),
\]
or
\[
P(A\mid B)>P(A).
\]

Similarly, since \(P(A)>0\),
\[
P(A\cap B)>P(A)P(B)
\]
is equivalent to
\[
\frac{P(A\cap B)}{P(A)}>P(B),
\]
which is
\[
P(B\mid A)>P(B).
\]

Hence,
\[
\operatorname{Cov}(I_A,I_B)>0
\]
if and only if
\[
P(A\mid B)>P(A),
\]
and if and only if
\[
P(B\mid A)>P(B).
\]

\end{solution}


%%%%%%%%%%%%%%%%%%%%%%%%%%%%%%%%%%%%%%%%%%%%%%%%%%%%%%%%%%%%%%
%%%%%%%%%%%%%%%%%%%%%%  PROBLEM 35  %%%%%%%%%%%%%%%%%%%%%%%%%%
%%%%%%%%%%%%%%%%%%%%%%%%%%%%%%%%%%%%%%%%%%%%%%%%%%%%%%%%%%%%%%

\newpage
\begin{problem}

Let \(X\) and \(Y\) be integrable random variables.

Prove the \textbf{Law of Total Expectation}, also called the
\textbf{Tower Property}:
\[
E\!\left[E(X\mid Y)\right]
=
E(X).
\]

Give a proof for the jointly continuous case.

\end{problem}

\begin{solution}

Suppose \(X\) and \(Y\) have joint density \(f_{X,Y}(x,y)\).

By definition,
\[
E(X\mid Y=y)
=
\int_{-\infty}^{\infty}
x f_{X\mid Y}(x\mid y)\,dx.
\]

Therefore,
\[
\begin{aligned}
E[E(X\mid Y)]
&=
\int_{-\infty}^{\infty}
E(X\mid Y=y)f_Y(y)\,dy\\
&=
\int_{-\infty}^{\infty}
\left[
\int_{-\infty}^{\infty}
x f_{X\mid Y}(x\mid y)\,dx
\right]
f_Y(y)\,dy.
\end{aligned}
\]

Interchanging the order of integration,
\[
\begin{aligned}
E[E(X\mid Y)]
&=
\int_{-\infty}^{\infty}
x
\left[
\int_{-\infty}^{\infty}
f_{X\mid Y}(x\mid y)f_Y(y)\,dy
\right]dx.
\end{aligned}
\]

Since
\[
f_{X\mid Y}(x\mid y)
=
\frac{f_{X,Y}(x,y)}{f_Y(y)},
\]
we have
\[
f_{X\mid Y}(x\mid y)f_Y(y)
=
f_{X,Y}(x,y).
\]

Hence,
\[
\begin{aligned}
E[E(X\mid Y)]
&=
\int_{-\infty}^{\infty}
x
\left[
\int_{-\infty}^{\infty}
f_{X,Y}(x,y)\,dy
\right]dx\\
&=
\int_{-\infty}^{\infty}
x f_X(x)\,dx\\
&=
E(X).
\end{aligned}
\]

Therefore,
\[
E[E(X\mid Y)]=E(X).
\]

\end{solution}


%%%%%%%%%%%%%%%%%%%%%%%%%%%%%%%%%%%%%%%%%%%%%%%%%%%%%%%%%%%%%%
%%%%%%%%%%%%%%%%%%%%%%  PROBLEM 36  %%%%%%%%%%%%%%%%%%%%%%%%%%
%%%%%%%%%%%%%%%%%%%%%%%%%%%%%%%%%%%%%%%%%%%%%%%%%%%%%%%%%%%%%%

\newpage
\begin{problem}

Let \(X\) and \(Y\) be random variables with \(E(X^2)<\infty\).

Prove the \textbf{Law of Total Variance}:
\[
\operatorname{Var}(X)
=
E\!\left[\operatorname{Var}(X\mid Y)\right]
+
\operatorname{Var}\!\left(E(X\mid Y)\right).
\]

\end{problem}

\begin{solution}

Recall that
\[
\operatorname{Var}(X)
=
E(X^2)-[E(X)]^2.
\]

For each value of \(Y\),
\[
\operatorname{Var}(X\mid Y)
=
E(X^2\mid Y)-[E(X\mid Y)]^2.
\]

Taking expectations,
\[
\begin{aligned}
E[\operatorname{Var}(X\mid Y)]
&=
E[E(X^2\mid Y)]
-
E\left([E(X\mid Y)]^2\right).
\end{aligned}
\]

By the Law of Total Expectation,
\[
E[E(X^2\mid Y)]
=
E(X^2).
\]

Therefore,
\[
E[\operatorname{Var}(X\mid Y)]
=
E(X^2)
-
E\left([E(X\mid Y)]^2\right).
\]

Also,
\[
\begin{aligned}
\operatorname{Var}(E(X\mid Y))
&=
E\left([E(X\mid Y)]^2\right)
-
\left(E[E(X\mid Y)]\right)^2.
\end{aligned}
\]

Again using the Law of Total Expectation,
\[
E[E(X\mid Y)]
=
E(X).
\]

Thus,
\[
\operatorname{Var}(E(X\mid Y))
=
E\left([E(X\mid Y)]^2\right)
-
[E(X)]^2.
\]

Adding the two expressions,
\[
\begin{aligned}
&E[\operatorname{Var}(X\mid Y)]
+
\operatorname{Var}(E(X\mid Y))\\
&=
E(X^2)
-
E\left([E(X\mid Y)]^2\right)
+
E\left([E(X\mid Y)]^2\right)
-
[E(X)]^2\\
&=
E(X^2)-[E(X)]^2\\
&=
\operatorname{Var}(X).
\end{aligned}
\]

Therefore,
\[
\operatorname{Var}(X)
=
E[\operatorname{Var}(X\mid Y)]
+
\operatorname{Var}(E(X\mid Y)).
\]

\end{solution}


%%%%%%%%%%%%%%%%%%%%%%%%%%%%%%%%%%%%%%%%%%%%%%%%%%%%%%%%%%%%%%
%%%%%%%%%%%%%%%%%%%%%%  PROBLEM 37  %%%%%%%%%%%%%%%%%%%%%%%%%%
%%%%%%%%%%%%%%%%%%%%%%%%%%%%%%%%%%%%%%%%%%%%%%%%%%%%%%%%%%%%%%

\newpage
\begin{problem}

Let \(X\) and \(Y\) be continuous random variables with joint probability
density function
\[
f(x,y)
=
\begin{cases}
\dfrac32(x^2+y^2),
&
0<x<1,\quad 0<y<1,\\[2mm]
0,
&
\text{otherwise}.
\end{cases}
\]

Find the random variable
\[
E(X\mid Y).
\]

\end{problem}

\begin{solution}

We first find the marginal density of \(Y\).

For \(0<y<1\),
\[
\begin{aligned}
f_Y(y)
&=
\int_0^1
\frac32(x^2+y^2)\,dx\\
&=
\frac32
\left(
\frac13+y^2
\right)\\
&=
\frac12+\frac32y^2\\
&=
\frac{3y^2+1}{2}.
\end{aligned}
\]

Thus,
\[
\begin{aligned}
f_{X\mid Y}(x\mid y)
&=
\frac{f(x,y)}{f_Y(y)}\\
&=
\frac{
\frac32(x^2+y^2)
}{
\frac12(3y^2+1)
}\\
&=
\frac{3(x^2+y^2)}
{3y^2+1},
\end{aligned}
\]
for \(0<x<1\).

Therefore,
\[
\begin{aligned}
E(X\mid Y=y)
&=
\int_0^1
x f_{X\mid Y}(x\mid y)\,dx\\
&=
\frac{3}{3y^2+1}
\int_0^1
(x^3+xy^2)\,dx\\
&=
\frac{3}{3y^2+1}
\left(
\frac14+\frac{y^2}{2}
\right)\\
&=
\frac{3(2y^2+1)}
{4(3y^2+1)}.
\end{aligned}
\]

Hence the conditional expectation as a random variable is
\[
E(X\mid Y)
=
\frac{3(2Y^2+1)}
{4(3Y^2+1)}.
\]

\end{solution}


%%%%%%%%%%%%%%%%%%%%%%%%%%%%%%%%%%%%%%%%%%%%%%%%%%%%%%%%%%%%%%
%%%%%%%%%%%%%%%%%%%%%%  PROBLEM 38  %%%%%%%%%%%%%%%%%%%%%%%%%%
%%%%%%%%%%%%%%%%%%%%%%%%%%%%%%%%%%%%%%%%%%%%%%%%%%%%%%%%%%%%%%

\newpage
\begin{problem}

For given random variables \(Y\) and \(Z\), suppose that \(X\) is selected
according to
\[
X=
\begin{cases}
Y, & \text{with probability }p,\\
Z, & \text{with probability }1-p,
\end{cases}
\]
where the random choice determining whether \(Y\) or \(Z\) is selected is
independent of \(Y\) and \(Z\).

Find \(E(X)\) in terms of \(E(Y)\) and \(E(Z)\).

\end{problem}

\begin{solution}

Let \(I\) be the indicator of the event that \(Y\) is selected. Then
\[
P(I=1)=p,
\qquad
P(I=0)=1-p,
\]
and
\[
X=IY+(1-I)Z.
\]

Since \(I\) is independent of \(Y\) and \(Z\),
\[
E(IY)=E(I)E(Y)=pE(Y),
\]
and
\[
E[(1-I)Z]
=
E(1-I)E(Z)
=
(1-p)E(Z).
\]

Therefore,
\[
\begin{aligned}
E(X)
&=
E[IY+(1-I)Z]\\
&=
pE(Y)+(1-p)E(Z).
\end{aligned}
\]

Hence,
\[
E(X)
=
pE(Y)+(1-p)E(Z).
\]

\end{solution}





%%%%%%%%%%%%%%%%%%%%%%%%%%%%%%%%%%%%%%%%%%%%%%%%%%%%%%%%%%%%%%
%%%%%%%%%%%%%%%%%%%%%%  PROBLEM 39  %%%%%%%%%%%%%%%%%%%%%%%%%%
%%%%%%%%%%%%%%%%%%%%%%%%%%%%%%%%%%%%%%%%%%%%%%%%%%%%%%%%%%%%%%

\newpage
\begin{problem}

Let \(X\) be a continuous random variable with probability density function
\[
f(x)
=
\begin{cases}
6x(1-x), & 0\leq x\leq1,\\
0, & \text{otherwise}.
\end{cases}
\]

\begin{enumerate}[label=(\alph*)]
    \item Find the moment-generating function \(M_X(t)\).
    \item Using \(M_X(t)\), find \(E(X)\).
\end{enumerate}

\end{problem}

\begin{solution}

\begin{enumerate}[label=(\alph*)]

\item By definition,
\[
M_X(t)
=
E(e^{tX})
=
\int_{-\infty}^{\infty}
e^{tx}f(x)\,dx.
\]

Therefore,
\[
M_X(t)
=
6\int_0^1
x(1-x)e^{tx}\,dx.
\]

Writing
\[
x(1-x)=x-x^2,
\]
we have
\[
M_X(t)
=
6\left[
\int_0^1 xe^{tx}\,dx
-
\int_0^1 x^2e^{tx}\,dx
\right].
\]

After integration by parts,
\[
M_X(t)
=
\frac{
6\left[
t+(t-2)e^t+2
\right]
}{t^3},
\qquad t\neq0.
\]

At \(t=0\),
\[
M_X(0)=1.
\]

Thus,
\[
M_X(t)
=
\begin{cases}
\dfrac{
6\left[t+(t-2)e^t+2\right]
}{t^3},
&
t\neq0,\\[3mm]
1,
&
t=0.
\end{cases}
\]

\item Recall that
\[
E(X)=M_X'(0).
\]

To evaluate this derivative conveniently, expand \(M_X(t)\) around \(t=0\).
Using
\[
e^t
=
1+t+\frac{t^2}{2}
+\frac{t^3}{6}
+\frac{t^4}{24}
+\cdots,
\]
we obtain
\[
M_X(t)
=
1+\frac12t+\frac{3}{20}t^2+\cdots.
\]

Therefore,
\[
M_X'(0)=\frac12.
\]

Hence,
\[
E(X)=\frac12.
\]

\end{enumerate}

\end{solution}


%%%%%%%%%%%%%%%%%%%%%%%%%%%%%%%%%%%%%%%%%%%%%%%%%%%%%%%%%%%%%%
%%%%%%%%%%%%%%%%%%%%%%  PROBLEM 40  %%%%%%%%%%%%%%%%%%%%%%%%%%
%%%%%%%%%%%%%%%%%%%%%%%%%%%%%%%%%%%%%%%%%%%%%%%%%%%%%%%%%%%%%%

\newpage
\begin{problem}

Let \(X\) be uniformly distributed on the interval \((a,b)\), where
\(a<b\). Find the moment-generating function of \(X\).

\end{problem}

\begin{solution}

Since
\[
X\sim\operatorname{Unif}(a,b),
\]
its density is
\[
f_X(x)
=
\frac1{b-a},
\qquad a<x<b.
\]

Therefore,
\[
\begin{aligned}
M_X(t)
&=
E(e^{tX})\\
&=
\frac1{b-a}
\int_a^b e^{tx}\,dx.
\end{aligned}
\]

For \(t\neq0\),
\[
\begin{aligned}
M_X(t)
&=
\frac1{b-a}
\left[
\frac{e^{tx}}{t}
\right]_a^b\\
&=
\frac{e^{bt}-e^{at}}
{(b-a)t}.
\end{aligned}
\]

Also,
\[
M_X(0)=1.
\]

Thus,
\[
M_X(t)
=
\begin{cases}
\dfrac{e^{bt}-e^{at}}
{(b-a)t},
&
t\neq0,\\[3mm]
1,
&
t=0.
\end{cases}
\]

\end{solution}


%%%%%%%%%%%%%%%%%%%%%%%%%%%%%%%%%%%%%%%%%%%%%%%%%%%%%%%%%%%%%%
%%%%%%%%%%%%%%%%%%%%%%  PROBLEM 41  %%%%%%%%%%%%%%%%%%%%%%%%%%
%%%%%%%%%%%%%%%%%%%%%%%%%%%%%%%%%%%%%%%%%%%%%%%%%%%%%%%%%%%%%%

\newpage
\begin{problem}

Let \(X\) be a geometric random variable with parameter \(p\), using the
convention
\[
P(X=k)=p(1-p)^{k-1},
\qquad
k=1,2,\ldots.
\]

Let
\[
q=1-p.
\]

\begin{enumerate}[label=(\alph*)]
    \item Show that
    \[
    M_X(t)
    =
    \frac{pe^t}{1-qe^t},
    \qquad
    t<-\ln q.
    \]

    \item Use \(M_X(t)\) to find \(E(X)\) and
    \(\operatorname{Var}(X)\).
\end{enumerate}

\end{problem}

\begin{solution}

\begin{enumerate}[label=(\alph*)]

\item By definition,
\[
\begin{aligned}
M_X(t)
&=
E(e^{tX})\\
&=
\sum_{k=1}^{\infty}
e^{tk}p q^{k-1}.
\end{aligned}
\]

Factor out \(pe^t\):
\[
M_X(t)
=
pe^t
\sum_{k=1}^{\infty}
(qe^t)^{k-1}.
\]

The geometric series converges when
\[
|qe^t|<1.
\]

Since \(q>0\), this is equivalent to
\[
t<-\ln q.
\]

Therefore,
\[
M_X(t)
=
\frac{pe^t}{1-qe^t},
\qquad
t<-\ln q.
\]

\item Differentiate:
\[
M_X'(t)
=
\frac{pe^t}
{(1-qe^t)^2}.
\]

Hence,
\[
E(X)
=
M_X'(0)
=
\frac{p}{(1-q)^2}.
\]

Since
\[
1-q=p,
\]
we obtain
\[
E(X)=\frac1p.
\]

Differentiate again:
\[
M_X''(t)
=
\frac{
pe^t(1+qe^t)
}{
(1-qe^t)^3
}.
\]

Therefore,
\[
E(X^2)
=
M_X''(0)
=
\frac{1+q}{p^2}.
\]

Thus,
\[
\begin{aligned}
\operatorname{Var}(X)
&=
E(X^2)-[E(X)]^2\\
&=
\frac{1+q}{p^2}
-
\frac1{p^2}\\
&=
\frac{q}{p^2}.
\end{aligned}
\]

Since \(q=1-p\),
\[
\operatorname{Var}(X)
=
\frac{1-p}{p^2}.
\]

\end{enumerate}

\end{solution}


%%%%%%%%%%%%%%%%%%%%%%%%%%%%%%%%%%%%%%%%%%%%%%%%%%%%%%%%%%%%%%
%%%%%%%%%%%%%%%%%%%%%%  PROBLEM 42  %%%%%%%%%%%%%%%%%%%%%%%%%%
%%%%%%%%%%%%%%%%%%%%%%%%%%%%%%%%%%%%%%%%%%%%%%%%%%%%%%%%%%%%%%

\newpage
\begin{problem}

Suppose that
\[
M_X(t)=\frac1{1-t},
\qquad t<1,
\]
is the moment-generating function of a random variable \(X\).

Find the moment-generating function of
\[
Y=2X+1.
\]

\end{problem}

\begin{solution}

By definition,
\[
M_Y(t)
=
E(e^{tY}).
\]

Since
\[
Y=2X+1,
\]
we have
\[
\begin{aligned}
M_Y(t)
&=
E\left(e^{t(2X+1)}\right)\\
&=
e^tE(e^{2tX})\\
&=
e^tM_X(2t).
\end{aligned}
\]

Therefore,
\[
M_Y(t)
=
e^t
\frac1{1-2t}.
\]

The MGF \(M_X(2t)\) is finite provided
\[
2t<1,
\]
that is,
\[
t<\frac12.
\]

Hence,
\[
M_Y(t)
=
\frac{e^t}{1-2t},
\qquad
t<\frac12.
\]

\end{solution}


%%%%%%%%%%%%%%%%%%%%%%%%%%%%%%%%%%%%%%%%%%%%%%%%%%%%%%%%%%%%%%
%%%%%%%%%%%%%%%%%%%%%%  PROBLEM 43  %%%%%%%%%%%%%%%%%%%%%%%%%%
%%%%%%%%%%%%%%%%%%%%%%%%%%%%%%%%%%%%%%%%%%%%%%%%%%%%%%%%%%%%%%

\newpage
\begin{problem}

For a random variable \(X\),
\[
M_X(t)
=
\left(
\frac{2}{2-t}
\right)^3.
\]

Find \(E(X)\) and \(\operatorname{Var}(X)\).

\end{problem}

\begin{solution}

Write
\[
M_X(t)
=
8(2-t)^{-3}.
\]

Differentiating,
\[
M_X'(t)
=
24(2-t)^{-4}.
\]

Therefore,
\[
E(X)
=
M_X'(0)
=
\frac{24}{16}
=
\frac32.
\]

Differentiating again,
\[
M_X''(t)
=
96(2-t)^{-5}.
\]

Thus,
\[
E(X^2)
=
M_X''(0)
=
\frac{96}{32}
=
3.
\]

Hence,
\[
\begin{aligned}
\operatorname{Var}(X)
&=
E(X^2)-[E(X)]^2\\
&=
3-\left(\frac32\right)^2\\
&=
3-\frac94\\
&=
\frac34.
\end{aligned}
\]

Therefore,
\[
E(X)=\frac32,
\qquad
\operatorname{Var}(X)=\frac34.
\]

\end{solution}


%%%%%%%%%%%%%%%%%%%%%%%%%%%%%%%%%%%%%%%%%%%%%%%%%%%%%%%%%%%%%%
%%%%%%%%%%%%%%%%%%%%%%  PROBLEM 44  %%%%%%%%%%%%%%%%%%%%%%%%%%
%%%%%%%%%%%%%%%%%%%%%%%%%%%%%%%%%%%%%%%%%%%%%%%%%%%%%%%%%%%%%%

\newpage
\begin{problem}

Let
\[
Z\sim N(0,1).
\]

Using
\[
M_Z(t)=e^{t^2/2},
\]
calculate
\[
E(Z^n),
\]
where \(n\) is a positive integer.

\end{problem}

\begin{solution}

The moment-generating function has the power-series representation
\[
M_Z(t)
=
\sum_{n=0}^{\infty}
\frac{E(Z^n)}{n!}t^n.
\]

On the other hand,
\[
M_Z(t)
=
e^{t^2/2}.
\]

Expanding the exponential,
\[
\begin{aligned}
e^{t^2/2}
&=
\sum_{k=0}^{\infty}
\frac1{k!}
\left(\frac{t^2}{2}\right)^k\\
&=
\sum_{k=0}^{\infty}
\frac{t^{2k}}{2^k k!}.
\end{aligned}
\]

There are no odd powers of \(t\). Therefore,
\[
E(Z^n)=0
\]
whenever \(n\) is odd.

If
\[
n=2k,
\]
then comparison of the coefficient of \(t^{2k}\) gives
\[
\frac{E(Z^{2k})}{(2k)!}
=
\frac1{2^k k!}.
\]

Thus,
\[
E(Z^{2k})
=
\frac{(2k)!}{2^k k!}.
\]

Equivalently,
\[
E(Z^{2k})
=
(2k-1)!!.
\]

Therefore,
\[
E(Z^n)
=
\begin{cases}
0,
&
n\text{ odd},\\[2mm]
\dfrac{n!}
{2^{n/2}(n/2)!},
&
n\text{ even}.
\end{cases}
\]

\end{solution}


%%%%%%%%%%%%%%%%%%%%%%%%%%%%%%%%%%%%%%%%%%%%%%%%%%%%%%%%%%%%%%
%%%%%%%%%%%%%%%%%%%%%%  PROBLEM 45  %%%%%%%%%%%%%%%%%%%%%%%%%%
%%%%%%%%%%%%%%%%%%%%%%%%%%%%%%%%%%%%%%%%%%%%%%%%%%%%%%%%%%%%%%

\newpage
\begin{problem}

Let
\[
X_1,\ldots,X_n
\]
be independent random variables with
\[
X_i\sim N(\mu_i,\sigma_i^2),
\qquad
i=1,\ldots,n.
\]

Let
\[
a_1,\ldots,a_n
\]
be constants.

Using moment-generating functions, prove that
\[
\sum_{i=1}^{n}a_iX_i
\sim
N\left(
\sum_{i=1}^{n}a_i\mu_i,
\,
\sum_{i=1}^{n}a_i^2\sigma_i^2
\right).
\]

Deduce that if
\[
X_1,\ldots,X_n
\overset{\text{i.i.d.}}{\sim}
N(\mu,\sigma^2),
\]
then
\[
\overline X
=
\frac1n\sum_{i=1}^{n}X_i
\sim
N\left(
\mu,\frac{\sigma^2}{n}
\right).
\]

\end{problem}

\begin{solution}

For
\[
X_i\sim N(\mu_i,\sigma_i^2),
\]
the MGF is
\[
M_{X_i}(t)
=
\exp\left(
\mu_i t+\frac12\sigma_i^2t^2
\right).
\]

Let
\[
S=\sum_{i=1}^{n}a_iX_i.
\]

Then
\[
M_S(t)
=
E\left[
e^{t\sum_{i=1}^{n}a_iX_i}
\right].
\]

Since the \(X_i\)'s are independent,
\[
\begin{aligned}
M_S(t)
&=
\prod_{i=1}^{n}
E(e^{ta_iX_i})\\
&=
\prod_{i=1}^{n}
M_{X_i}(a_it).
\end{aligned}
\]

Therefore,
\[
\begin{aligned}
M_S(t)
&=
\prod_{i=1}^{n}
\exp\left(
a_i\mu_i t
+
\frac12a_i^2\sigma_i^2t^2
\right)\\
&=
\exp\left[
t\sum_{i=1}^{n}a_i\mu_i
+
\frac{t^2}{2}
\sum_{i=1}^{n}a_i^2\sigma_i^2
\right].
\end{aligned}
\]

This is the MGF of a Normal random variable with mean
\[
\sum_{i=1}^{n}a_i\mu_i
\]
and variance
\[
\sum_{i=1}^{n}a_i^2\sigma_i^2.
\]

Hence,
\[
\sum_{i=1}^{n}a_iX_i
\sim
N\left(
\sum_{i=1}^{n}a_i\mu_i,
\,
\sum_{i=1}^{n}a_i^2\sigma_i^2
\right).
\]

For the sample mean, take
\[
a_i=\frac1n,
\qquad
\mu_i=\mu,
\qquad
\sigma_i^2=\sigma^2.
\]

Then
\[
E(\overline X)
=
\sum_{i=1}^{n}
\frac1n\mu
=
\mu,
\]
and
\[
\operatorname{Var}(\overline X)
=
\sum_{i=1}^{n}
\frac1{n^2}\sigma^2
=
\frac{\sigma^2}{n}.
\]

Therefore,
\[
\overline X
\sim
N\left(
\mu,\frac{\sigma^2}{n}
\right).
\]

\end{solution}


%%%%%%%%%%%%%%%%%%%%%%%%%%%%%%%%%%%%%%%%%%%%%%%%%%%%%%%%%%%%%%
%%%%%%%%%%%%%%%%%%%%%%  PROBLEM 46  %%%%%%%%%%%%%%%%%%%%%%%%%%
%%%%%%%%%%%%%%%%%%%%%%%%%%%%%%%%%%%%%%%%%%%%%%%%%%%%%%%%%%%%%%

\newpage
\begin{problem}

Let
\[
X_1,X_2,\ldots,X_n
\]
be independent geometric random variables, each with parameter \(p\), using
the convention
\[
P(X_i=k)
=
p(1-p)^{k-1},
\qquad k=1,2,\ldots.
\]

Using moment-generating functions, prove that
\[
X_1+\cdots+X_n
\]
has a negative binomial distribution with parameters \((n,p)\).

\end{problem}

\begin{solution}

Let
\[
q=1-p.
\]

The MGF of a geometric random variable under this convention is
\[
M_{X_i}(t)
=
\frac{pe^t}{1-qe^t},
\qquad
t<-\ln q.
\]

Define
\[
S_n=X_1+\cdots+X_n.
\]

Since \(X_1,\ldots,X_n\) are independent,
\[
\begin{aligned}
M_{S_n}(t)
&=
\prod_{i=1}^{n}M_{X_i}(t)\\
&=
\left(
\frac{pe^t}{1-qe^t}
\right)^n.
\end{aligned}
\]

Now recall that if
\[
W\sim\operatorname{NegBin}(n,p),
\]
where \(W\) denotes the trial on which the \(n\)-th success occurs, then
\[
P(W=k)
=
\binom{k-1}{n-1}
p^nq^{k-n},
\qquad
k=n,n+1,\ldots.
\]

Its MGF is
\[
M_W(t)
=
\left(
\frac{pe^t}{1-qe^t}
\right)^n.
\]

Thus,
\[
M_{S_n}(t)=M_W(t)
\]
in a neighborhood of \(0\).

By uniqueness of moment-generating functions,
\[
S_n
=
X_1+\cdots+X_n
\sim
\operatorname{NegBin}(n,p).
\]

\end{solution}


%%%%%%%%%%%%%%%%%%%%%%%%%%%%%%%%%%%%%%%%%%%%%%%%%%%%%%%%%%%%%%
%%%%%%%%%%%%%%%%%%%%%%  PROBLEM 47  %%%%%%%%%%%%%%%%%%%%%%%%%%
%%%%%%%%%%%%%%%%%%%%%%%%%%%%%%%%%%%%%%%%%%%%%%%%%%%%%%%%%%%%%%

\newpage
\begin{problem}

Let \(X\) and \(Y\) be independent binomial random variables with
\[
X\sim\operatorname{Bin}(n,p),
\qquad
Y\sim\operatorname{Bin}(m,p).
\]

Calculate
\[
P(X=i\mid X+Y=j)
\]
and interpret the result.

\end{problem}

\begin{solution}

Since \(X\) and \(Y\) are independent,
\[
X+Y
\sim
\operatorname{Bin}(n+m,p).
\]

Therefore,
\[
P(X=i\mid X+Y=j)
=
\frac{
P(X=i,Y=j-i)
}{
P(X+Y=j)
}.
\]

The numerator is
\[
\begin{aligned}
P(X=i,Y=j-i)
&=
P(X=i)P(Y=j-i)\\
&=
\binom ni
p^i(1-p)^{n-i}\\
&\qquad\times
\binom m{j-i}
p^{j-i}(1-p)^{m-j+i}.
\end{aligned}
\]

Thus,
\[
P(X=i,Y=j-i)
=
\binom ni
\binom m{j-i}
p^j(1-p)^{n+m-j}.
\]

Also,
\[
P(X+Y=j)
=
\binom{n+m}{j}
p^j(1-p)^{n+m-j}.
\]

Hence,
\[
P(X=i\mid X+Y=j)
=
\frac{
\binom ni
\binom m{j-i}
}{
\binom{n+m}{j}
},
\]
for
\[
\max\{0,j-m\}
\leq i
\leq
\min\{n,j\}.
\]

This is the probability mass function of a hypergeometric random variable.

Conditional on there being exactly \(j\) total successes among the
\(n+m\) Bernoulli trials, all choices of the \(j\) successful trials are
equally likely. The random variable \(X\) counts how many of those \(j\)
successes occur among the first \(n\) trials.

Therefore,
\[
X\mid(X+Y=j)
\sim
\operatorname{Hypergeom}(n+m,n,j).
\]

\end{solution}


%%%%%%%%%%%%%%%%%%%%%%%%%%%%%%%%%%%%%%%%%%%%%%%%%%%%%%%%%%%%%%
%%%%%%%%%%%%%%%%%%%%%%  PROBLEM 48  %%%%%%%%%%%%%%%%%%%%%%%%%%
%%%%%%%%%%%%%%%%%%%%%%%%%%%%%%%%%%%%%%%%%%%%%%%%%%%%%%%%%%%%%%

\newpage
\begin{problem}

Let \(X,Y,Z\) be independent Poisson random variables with parameters
\[
\lambda_1,\lambda_2,\lambda_3,
\]
respectively.

For
\[
y=0,1,\ldots,t,
\]
calculate
\[
P(Y=y\mid X+Y+Z=t).
\]

\end{problem}

\begin{solution}

Let
\[
T=X+Y+Z.
\]

Since independent Poisson random variables add,
\[
T
\sim
\operatorname{Pois}
(\lambda_1+\lambda_2+\lambda_3).
\]

Also,
\[
X+Z
\sim
\operatorname{Pois}(\lambda_1+\lambda_3),
\]
and \(X+Z\) is independent of \(Y\).

Therefore,
\[
\begin{aligned}
P(Y=y\mid T=t)
&=
\frac{
P(Y=y,\ X+Z=t-y)
}{
P(T=t)
}\\
&=
\frac{
P(Y=y)
P(X+Z=t-y)
}{
P(T=t)
}.
\end{aligned}
\]

Substituting the Poisson probabilities,
\[
\begin{aligned}
P(Y=y\mid T=t)
&=
\frac{
e^{-\lambda_2}
\dfrac{\lambda_2^y}{y!}
\,
e^{-(\lambda_1+\lambda_3)}
\dfrac{(\lambda_1+\lambda_3)^{t-y}}{(t-y)!}
}{
e^{-(\lambda_1+\lambda_2+\lambda_3)}
\dfrac{
(\lambda_1+\lambda_2+\lambda_3)^t
}{t!}
}.
\end{aligned}
\]

The exponential terms cancel, giving
\[
P(Y=y\mid T=t)
=
\binom ty
\frac{
\lambda_2^y
(\lambda_1+\lambda_3)^{t-y}
}{
(\lambda_1+\lambda_2+\lambda_3)^t
}.
\]

Equivalently,
\[
P(Y=y\mid T=t)
=
\binom ty
\left(
\frac{\lambda_2}
{\lambda_1+\lambda_2+\lambda_3}
\right)^y
\left(
\frac{\lambda_1+\lambda_3}
{\lambda_1+\lambda_2+\lambda_3}
\right)^{t-y}.
\]

Thus,
\[
Y\mid(X+Y+Z=t)
\sim
\operatorname{Bin}
\left(
t,
\frac{\lambda_2}
{\lambda_1+\lambda_2+\lambda_3}
\right).
\]

\end{solution}


%%%%%%%%%%%%%%%%%%%%%%%%%%%%%%%%%%%%%%%%%%%%%%%%%%%%%%%%%%%%%%
%%%%%%%%%%%%%%%%%%%%%%  PROBLEM 49  %%%%%%%%%%%%%%%%%%%%%%%%%%
%%%%%%%%%%%%%%%%%%%%%%%%%%%%%%%%%%%%%%%%%%%%%%%%%%%%%%%%%%%%%%

\newpage%%%%%%%%%%%%%%%%%%%%%%%%%%%%%%%%%%%%%%%%%%%%%%%%%%%%%%%%%%%%%%
%%%%%%%%%%%%%%%%%%%%  INEQUALITIES PROBLEM  %%%%%%%%%%%%%%%%%%
%%%%%%%%%%%%%%%%%%%%%%%%%%%%%%%%%%%%%%%%%%%%%%%%%%%%%%%%%%%%%%

\newpage
\begin{problem}[]

Let \(X\) and \(Y\) be i.i.d. positive random variables, and let \(c>0\).

For each part below, fill in the appropriate equality or inequality symbol.

Write \(=\) if the two sides are always equal, \(\leq\) if the left-hand
side is less than or equal to the right-hand side but not necessarily equal,
and similarly for \(\geq\). If no relation holds in general, write \(?\).

\begin{enumerate}[label=(\alph*)]

    \item
    \[
    E(\log X)
    \ \underline{\hspace{1cm}}\
    \log(E X).
    \]

    \item
    \[
    E(X)
    \ \underline{\hspace{1cm}}\
    \sqrt{E(X^2)}.
    \]

    \item
    \[
    E(\sin^2 X)+E(\cos^2 X)
    \ \underline{\hspace{1cm}}\
    1.
    \]

    \item
    \[
    E(|X|)
    \ \underline{\hspace{1cm}}\
    \sqrt{E(X^2)}.
    \]

    \item
    \[
    P(X>c)
    \ \underline{\hspace{1cm}}\
    \frac{E(X^3)}{c^3}.
    \]

    \item
    \[
    P(X\leq Y)
    \ \underline{\hspace{1cm}}\
    P(X\geq Y).
    \]

    \item
    \[
    E(XY)
    \ \underline{\hspace{1cm}}\
    \sqrt{E(X^2)E(Y^2)}.
    \]

    \item
    \[
    P(X+Y>10)
    \ \underline{\hspace{1cm}}\
    P(X>5\text{ or }Y>5).
    \]

    \item
    \[
    E(\min(X,Y))
    \ \underline{\hspace{1cm}}\
    \min(E X,E Y).
    \]

    \item
    \[
    E\left(\frac{X}{Y}\right)
    \ \underline{\hspace{1cm}}\
    \frac{E X}{E Y}.
    \]

    \item
    \[
    E\left(X^2(X^2+1)\right)
    \ \underline{\hspace{1cm}}\
    E\left(X^2(Y^2+1)\right).
    \]

    \item
    \[
    E\left(\frac{X^3}{X^3+Y^3}\right)
    \ \underline{\hspace{1cm}}\
    E\left(\frac{Y^3}{X^3+Y^3}\right).
    \]

\end{enumerate}

\end{problem}


\begin{solution}

\begin{enumerate}[label=(\alph*)]

\item
\[
E(\log X)\leq \log(E X).
\]

Since \(\log x\) is concave on \((0,\infty)\), Jensen's inequality gives
\[
E(\log X)\leq \log(E X).
\]


\item
\[
E(X)\leq \sqrt{E(X^2)}.
\]

By Cauchy--Schwarz,
\[
[E(X)]^2\leq E(X^2).
\]

Since \(X>0\),
\[
E(X)\leq \sqrt{E(X^2)}.
\]


\item
\[
E(\sin^2 X)+E(\cos^2 X)=1.
\]

Indeed,
\[
\sin^2 X+\cos^2 X=1
\]
for every value of \(X\). Therefore,
\[
\begin{aligned}
E(\sin^2 X)+E(\cos^2 X)
&=
E(\sin^2 X+\cos^2 X)\\
&=
E(1)\\
&=
1.
\end{aligned}
\]


\item
\[
E(|X|)\leq \sqrt{E(X^2)}.
\]

By Cauchy--Schwarz,
\[
[E(|X|)]^2
\leq
E(|X|^2)
=
E(X^2).
\]

Hence,
\[
E(|X|)
\leq
\sqrt{E(X^2)}.
\]


\item
\[
P(X>c)\leq \frac{E(X^3)}{c^3}.
\]

Since \(X>0\),
\[
X>c
\quad\Longleftrightarrow\quad
X^3>c^3.
\]

Applying Markov's inequality to the nonnegative random variable \(X^3\),
\[
P(X>c)
=
P(X^3>c^3)
\leq
\frac{E(X^3)}{c^3}.
\]


\item
\[
P(X\leq Y)=P(X\geq Y).
\]

Since \(X\) and \(Y\) are identically distributed and independent, the joint
distribution of \((X,Y)\) is unchanged when \(X\) and \(Y\) are exchanged.

Therefore,
\[
P(X\leq Y)=P(Y\leq X)=P(X\geq Y).
\]


\item
\[
E(XY)\leq \sqrt{E(X^2)E(Y^2)}.
\]

This follows directly from the Cauchy--Schwarz inequality:
\[
|E(XY)|
\leq
\sqrt{E(X^2)E(Y^2)}.
\]

Since \(X\) and \(Y\) are positive,
\[
E(XY)\geq0,
\]
and hence
\[
E(XY)
\leq
\sqrt{E(X^2)E(Y^2)}.
\]


\item
\[
P(X+Y>10)
\leq
P(X>5\text{ or }Y>5).
\]

Indeed,
\[
\{X+Y>10\}
\subseteq
\{X>5\}\cup\{Y>5\}.
\]

Therefore,
\[
P(X+Y>10)
\leq
P(X>5\text{ or }Y>5).
\]


\item
\[
E(\min(X,Y))
\leq
\min(E X,E Y).
\]

Pointwise,
\[
\min(X,Y)\leq X
\]
and
\[
\min(X,Y)\leq Y.
\]

Taking expectations gives
\[
E(\min(X,Y))\leq E(X)
\]
and
\[
E(\min(X,Y))\leq E(Y).
\]

Therefore,
\[
E(\min(X,Y))
\leq
\min(E X,E Y).
\]


\item
\[
E\left(\frac{X}{Y}\right)
\ ?\
\frac{E X}{E Y}.
\]

No fixed relation holds in general.

By independence,
\[
E\left(\frac{X}{Y}\right)
=
E(X)E\left(\frac1Y\right),
\]
provided the expectations exist. This is generally not equal to
\[
\frac{E(X)}{E(Y)}.
\]


\item
\[
E\left(X^2(X^2+1)\right)
\geq
E\left(X^2(Y^2+1)\right).
\]

For the left-hand side,
\[
E\left(X^2(X^2+1)\right)
=
E(X^4)+E(X^2).
\]

For the right-hand side, independence gives
\[
\begin{aligned}
E\left(X^2(Y^2+1)\right)
&=
E(X^2Y^2)+E(X^2)\\
&=
E(X^2)E(Y^2)+E(X^2).
\end{aligned}
\]

Since \(X\) and \(Y\) are identically distributed,
\[
E(Y^2)=E(X^2),
\]
so
\[
E\left(X^2(Y^2+1)\right)
=
[E(X^2)]^2+E(X^2).
\]

Finally,
\[
E(X^4)\geq [E(X^2)]^2
\]
because
\[
\operatorname{Var}(X^2)\geq0.
\]

Therefore,
\[
E\left(X^2(X^2+1)\right)
\geq
E\left(X^2(Y^2+1)\right).
\]


\item
\[
E\left(\frac{X^3}{X^3+Y^3}\right)
=
E\left(\frac{Y^3}{X^3+Y^3}\right).
\]

Since \(X\) and \(Y\) are i.i.d., exchanging \(X\) and \(Y\) does not change
their joint distribution. Hence,
\[
E\left(\frac{X^3}{X^3+Y^3}\right)
=
E\left(\frac{Y^3}{X^3+Y^3}\right).
\]

In fact, since
\[
\frac{X^3}{X^3+Y^3}
+
\frac{Y^3}{X^3+Y^3}
=
1,
\]
the two equal expectations must each be
\[
\frac12.
\]

\end{enumerate}

\end{solution}


%%%%%%%%%%%%%%%%%%%%%%%%%%%%%%%%%%%%%%%%%%%%%%%%%%%%%%%%%%%%%%
%%%%%%%%%%%%%%%%%%%%%%  PROBLEM 50  %%%%%%%%%%%%%%%%%%%%%%%%%%
%%%%%%%%%%%%%%%%%%%%%%%%%%%%%%%%%%%%%%%%%%%%%%%%%%%%%%%%%%%%%%
%%%%%%%%%%%%%%%%%%%%%%%%%%%%%%%%%%%%%%%%%%%%%%%%%%%%%%%%%%%%%%
%%%%%%%%%%%%%%%  PROBABILITY INEQUALITIES PROBLEM  %%%%%%%%%%%
%%%%%%%%%%%%%%%%%%%%%%%%%%%%%%%%%%%%%%%%%%%%%%%%%%%%%%%%%%%%%%

\newpage
\begin{problem}

Let \(X\) be an exponential random variable with parameter \(1\), so that
\[
f_X(x)
=
\begin{cases}
e^{-x}, & x>0,\\
0, & \text{otherwise}.
\end{cases}
\]

Recall that
\[
E(X)=1,
\qquad
\operatorname{Var}(X)=1,
\]
and the moment generating function of \(X\) is
\[
M_X(t)=\frac{1}{1-t},
\qquad t<1.
\]

\begin{enumerate}[label=(\alph*)]

    \item Use Markov's inequality to find an upper bound for
    \[
    P(X\geq 10).
    \]

    \item Use Chebyshev's inequality to find an upper bound for
    \[
    P(X\geq 10).
    \]

    \item Recall that, for any \(t>0\),
    \[
    P(X\geq a)
    \leq
    e^{-ta}M_X(t).
    \]
    Use this Chernoff bound to show that
    \[
    P(X\geq 10)
    \leq
    \frac{e^{-10t}}{1-t},
    \qquad 0<t<1.
    \]

    Find the value of \(t\) that gives the best bound, and calculate the
    resulting Chernoff bound.

    \item Compare the three bounds obtained above. Which inequality gives
    the sharpest upper bound for \(P(X\geq10)\)?

    Also calculate the exact value of \(P(X\geq10)\) and compare it with
    the three bounds.

    \item Compare
    \[
    E\left(\frac{1}{1+X}\right)
    \qquad\text{and}\qquad
    \frac{1}{1+E(X)}
    \]
    i.e., determine whether \(=\), \(\leq\), or \(\geq\) always holds.

    \textit{Hint:} Consider the function
    \[
    g(x)=\frac{1}{1+x}.
    \]
    Check whether the conditions for Jensen's inequality are satisfied.

\end{enumerate}

\end{problem}


\begin{solution}

\begin{enumerate}[label=(\alph*)]

\item Markov's inequality states that, for a nonnegative random variable,
\[
P(X\geq a)
\leq
\frac{E(X)}{a}.
\]

Since
\[
E(X)=1,
\]
we obtain
\[
P(X\geq10)
\leq
\frac{1}{10}.
\]

Thus, the Markov bound is
\[
P(X\geq10)\leq 0.1.
\]


\item To apply Chebyshev's inequality, observe that
\[
X\geq10
\]
implies
\[
X-E(X)\geq9.
\]

Therefore,
\[
\{X\geq10\}
\subseteq
\{|X-E(X)|\geq9\}.
\]

Hence,
\[
\begin{aligned}
P(X\geq10)
&\leq
P(|X-E(X)|\geq9)\\
&\leq
\frac{\operatorname{Var}(X)}{9^2}\\
&=
\frac{1}{81}.
\end{aligned}
\]

Thus, the Chebyshev bound is
\[
P(X\geq10)
\leq
\frac{1}{81}
\approx
0.01235.
\]


\item For \(t>0\), the Chernoff bound gives
\[
P(X\geq10)
\leq
e^{-10t}M_X(t).
\]

Since
\[
M_X(t)=\frac{1}{1-t},
\qquad t<1,
\]
we have
\[
P(X\geq10)
\leq
\frac{e^{-10t}}{1-t},
\qquad 0<t<1.
\]

We now minimize
\[
h(t)=\frac{e^{-10t}}{1-t}.
\]

It is convenient to minimize its logarithm:
\[
\log h(t)
=
-10t-\log(1-t).
\]

Differentiating,
\[
\frac{d}{dt}\log h(t)
=
-10+\frac{1}{1-t}.
\]

Setting this equal to zero,
\[
-10+\frac{1}{1-t}=0.
\]

Thus,
\[
\frac{1}{1-t}=10,
\]
so
\[
1-t=\frac{1}{10},
\]
and therefore
\[
t=\frac{9}{10}.
\]

Substituting this value into the Chernoff bound,
\[
\begin{aligned}
P(X\geq10)
&\leq
\frac{e^{-10(9/10)}}{1-9/10}\\
&=
\frac{e^{-9}}{1/10}\\
&=
10e^{-9}.
\end{aligned}
\]

Thus, the optimized Chernoff bound is
\[
P(X\geq10)
\leq
10e^{-9}
\approx
0.001234.
\]


\item The three bounds are
\[
\begin{array}{c|c}
\text{Method} & \text{Upper bound}\\
\hline
\text{Markov} & \dfrac{1}{10}\\[2mm]
\text{Chebyshev} & \dfrac{1}{81}\\[2mm]
\text{Chernoff} & 10e^{-9}
\end{array}
\]

Numerically,
\[
\frac{1}{10}=0.1,
\]
\[
\frac{1}{81}\approx0.01235,
\]
and
\[
10e^{-9}\approx0.001234.
\]

Therefore,
\[
10e^{-9}
<
\frac{1}{81}
<
\frac{1}{10}.
\]

Hence, among these three inequalities, the Chernoff bound gives the
sharpest upper bound.

The exact probability can be calculated directly from the exponential
distribution:
\[
\begin{aligned}
P(X\geq10)
&=
\int_{10}^{\infty}e^{-x}\,dx\\
&=
e^{-10}.
\end{aligned}
\]

Numerically,
\[
e^{-10}\approx0.0000454.
\]

Thus,
\[
\underbrace{e^{-10}}_{\text{exact}}
<
\underbrace{10e^{-9}}_{\text{Chernoff}}
<
\underbrace{\frac{1}{81}}_{\text{Chebyshev}}
<
\underbrace{\frac{1}{10}}_{\text{Markov}}.
\]


\item Define
\[
g(x)=\frac{1}{1+x},
\qquad x>0.
\]

We have
\[
g'(x)
=
-\frac{1}{(1+x)^2}
\]
and
\[
g''(x)
=
\frac{2}{(1+x)^3}.
\]

Since
\[
g''(x)>0
\qquad\text{for }x>0,
\]
the function \(g\) is convex on the support of \(X\).

Therefore, Jensen's inequality gives
\[
g(E(X))
\leq
E(g(X)).
\]

Hence,
\[
\frac{1}{1+E(X)}
\leq
E\left(\frac{1}{1+X}\right).
\]

Equivalently,
\[
E\left(\frac{1}{1+X}\right)
\geq
\frac{1}{1+E(X)}.
\]

Since \(E(X)=1\),
\[
E\left(\frac{1}{1+X}\right)
\geq
\frac12.
\]

\end{enumerate}

\end{solution}


%%%%%%%%%%%%%%%%%%%%%%%%%%%%%%%%%%%%%%%%%%%%%%%%%%%%%%%%%%%%%%
%%%%%%%%%%%%%%%%%%%%%%  PROBLEM 51  %%%%%%%%%%%%%%%%%%%%%%%%%%
%%%%%%%%%%%%%%%%%%%%%%%%%%%%%%%%%%%%%%%%%%%%%%%%%%%%%%%%%%%%%%

\newpage
\begin{problem}

Suppose that the number \(X\) of letters handled by a post office on a given
day satisfies
\[
E(X)=10{,}000
\]
and
\[
\operatorname{Var}(X)=2000.
\]

Using Chebyshev's inequality, find a lower bound for
\[
P(8000<X<12{,}000).
\]

\end{problem}

\begin{solution}

We can write
\[
8000<X<12{,}000
\]
as
\[
-2000<X-10{,}000<2000.
\]

Thus,
\[
P(8000<X<12{,}000)
=
P(|X-10{,}000|<2000).
\]

Using the complement,
\[
P(|X-10{,}000|<2000)
=
1-
P(|X-10{,}000|\geq2000).
\]

By Chebyshev's inequality,
\[
P(|X-10{,}000|\geq2000)
\leq
\frac{2000}{2000^2}.
\]

Therefore,
\[
P(|X-10{,}000|\geq2000)
\leq
0.0005.
\]

Hence,
\[
P(8000<X<12{,}000)
\geq
1-0.0005
=
0.9995.
\]

Thus,
\[
P(8000<X<12{,}000)
\geq
0.9995.
\]

\end{solution}


%%%%%%%%%%%%%%%%%%%%%%%%%%%%%%%%%%%%%%%%%%%%%%%%%%%%%%%%%%%%%%
%%%%%%%%%%%%%%%%%%%%%%  PROBLEM 52  %%%%%%%%%%%%%%%%%%%%%%%%%%
%%%%%%%%%%%%%%%%%%%%%%%%%%%%%%%%%%%%%%%%%%%%%%%%%%%%%%%%%%%%%%

\newpage
\begin{problem}

Roll a fair six-sided die and let \(X\) denote the outcome.

\begin{enumerate}[label=(\alph*)]
    \item Calculate \(E(X)\) and \(\operatorname{Var}(X)\).

    \item Use Markov's inequality to obtain an upper bound for
    \[
    P(X\geq6).
    \]

    \item Use Chebyshev's inequality to obtain an upper bound for
    \[
    P\left(
    \left|X-\frac{21}{6}\right|
    \geq
    \frac32
    \right).
    \]

    \item Calculate both probabilities exactly and compare the exact
    probabilities with the bounds.
\end{enumerate}

\end{problem}

\begin{solution}

\begin{enumerate}[label=(\alph*)]

\item Since each outcome has probability \(1/6\),
\[
E(X)
=
\frac16(1+2+3+4+5+6)
=
\frac{21}{6}
=
\frac72.
\]

Also,
\[
E(X^2)
=
\frac16
(1^2+2^2+3^2+4^2+5^2+6^2)
=
\frac{91}{6}.
\]

Therefore,
\[
\begin{aligned}
\operatorname{Var}(X)
&=
E(X^2)-[E(X)]^2\\
&=
\frac{91}{6}
-
\frac{49}{4}\\
&=
\frac{35}{12}.
\end{aligned}
\]

\item By Markov's inequality,
\[
P(X\geq6)
\leq
\frac{E(X)}6.
\]

Hence,
\[
P(X\geq6)
\leq
\frac{7/2}{6}
=
\frac7{12}
\approx0.5833.
\]

\item By Chebyshev's inequality,
\[
P\left(
\left|X-\frac72\right|
\geq
\frac32
\right)
\leq
\frac{
35/12
}{
(3/2)^2
}.
\]

Therefore,
\[
P\left(
\left|X-\frac72\right|
\geq
\frac32
\right)
\leq
\frac{35}{27}
\approx1.296.
\]

Since probabilities cannot exceed \(1\), this bound is trivial. The useful
upper bound is simply
\[
1.
\]

\item Exactly,
\[
P(X\geq6)
=
P(X=6)
=
\frac16.
\]

Also,
\[
\left|X-\frac72\right|
\geq
\frac32
\]
occurs when
\[
X\leq2
\qquad\text{or}\qquad
X\geq5.
\]

Thus,
\[
P\left(
\left|X-\frac72\right|
\geq
\frac32
\right)
=
\frac46
=
\frac23.
\]

Therefore, both inequalities give valid upper bounds, but the bounds may be
far from the exact probabilities.

\end{enumerate}

\end{solution}


%%%%%%%%%%%%%%%%%%%%%%%%%%%%%%%%%%%%%%%%%%%%%%%%%%%%%%%%%%%%%%
%%%%%%%%%%%%%%%%%%%%%%  PROBLEM 53  %%%%%%%%%%%%%%%%%%%%%%%%%%
%%%%%%%%%%%%%%%%%%%%%%%%%%%%%%%%%%%%%%%%%%%%%%%%%%%%%%%%%%%%%%

\newpage
\begin{problem}

The scores on an achievement test have population mean
\[
\mu=500
\]
and population standard deviation
\[
\sigma=100.
\]

Let
\[
\overline X
\]
be the sample mean of a random sample of \(10\) students.

Without assuming that the population distribution is Normal, use Chebyshev's
inequality to find a lower bound for
\[
P(460<\overline X<540).
\]

\end{problem}

\begin{solution}

Since the sample consists of independent observations,
\[
E(\overline X)=500
\]
and
\[
\operatorname{Var}(\overline X)
=
\frac{\sigma^2}{n}
=
\frac{100^2}{10}
=
1000.
\]

Now,
\[
\begin{aligned}
P(460<\overline X<540)
&=
P(-40<\overline X-500<40)\\
&=
P(|\overline X-500|<40).
\end{aligned}
\]

Using the complement,
\[
P(|\overline X-500|<40)
=
1-
P(|\overline X-500|\geq40).
\]

By Chebyshev's inequality,
\[
P(|\overline X-500|\geq40)
\leq
\frac{1000}{40^2}.
\]

Thus,
\[
P(|\overline X-500|\geq40)
\leq
\frac{1000}{1600}
=
\frac58.
\]

Therefore,
\[
P(460<\overline X<540)
\geq
1-\frac58
=
\frac38.
\]

Hence,
\[
P(460<\overline X<540)
\geq
\frac38.
\]

\end{solution}


%%%%%%%%%%%%%%%%%%%%%%%%%%%%%%%%%%%%%%%%%%%%%%%%%%%%%%%%%%%%%%
%%%%%%%%%%%%%%%%%%%%%%  PROBLEM 54  %%%%%%%%%%%%%%%%%%%%%%%%%%
%%%%%%%%%%%%%%%%%%%%%%%%%%%%%%%%%%%%%%%%%%%%%%%%%%%%%%%%%%%%%%

\newpage
\begin{problem}

Suppose
\[
X_1,\ldots,X_n
\]
are independent Bernoulli random variables with unknown success probability
\(p\), and define
\[
\widehat p
=
\frac1n\sum_{i=1}^{n}X_i.
\]

Let
\[
\varepsilon>0
\qquad\text{and}\qquad
0<\alpha<1.
\]

Using Chebyshev's inequality, derive a sample-size condition that guarantees
\[
P(|\widehat p-p|<\varepsilon)
\geq
1-\alpha
\]
for every
\[
0<p<1.
\]

\end{problem}

\begin{solution}

Since
\[
X_i\sim\operatorname{Bernoulli}(p),
\]
we have
\[
E(X_i)=p
\]
and
\[
\operatorname{Var}(X_i)=p(1-p).
\]

Therefore,
\[
E(\widehat p)=p,
\]
and by independence,
\[
\operatorname{Var}(\widehat p)
=
\frac{p(1-p)}{n}.
\]

By Chebyshev's inequality,
\[
P(|\widehat p-p|\geq\varepsilon)
\leq
\frac{
p(1-p)
}{
n\varepsilon^2
}.
\]

Hence,
\[
P(|\widehat p-p|<\varepsilon)
\geq
1-
\frac{
p(1-p)
}{
n\varepsilon^2
}.
\]

To guarantee
\[
P(|\widehat p-p|<\varepsilon)
\geq
1-\alpha,
\]
it is sufficient that
\[
\frac{
p(1-p)
}{
n\varepsilon^2
}
\leq
\alpha.
\]

Equivalently,
\[
n
\geq
\frac{
p(1-p)
}{
\varepsilon^2\alpha
}.
\]

However, \(p\) is unknown. Since
\[
p(1-p)
\leq
\frac14
\]
for every \(0<p<1\), a sufficient condition that does not depend on \(p\) is
\[
n
\geq
\frac{
1
}{
4\varepsilon^2\alpha
}.
\]

Thus, choosing
\[
n
\geq
\frac1{4\varepsilon^2\alpha}
\]
guarantees
\[
P(|\widehat p-p|<\varepsilon)
\geq
1-\alpha
\]
for every Bernoulli success probability \(p\).

\end{solution}


%%%%%%%%%%%%%%%%%%%%%%%%%%%%%%%%%%%%%%%%%%%%%%%%%%%%%%%%%%%%%%
%%%%%%%%%%%%%%%%%%%%%%  PROBLEM 55  %%%%%%%%%%%%%%%%%%%%%%%%%%
%%%%%%%%%%%%%%%%%%%%%%%%%%%%%%%%%%%%%%%%%%%%%%%%%%%%%%%%%%%%%%

\newpage
\begin{problem}

A quality-control officer wants to estimate the proportion \(p\) of defective
nails produced by a company.

A random sample of \(n\) nails is inspected, and
\[
\widehat p
\]
denotes the fraction of defective nails in the sample.

Using Chebyshev's inequality, determine a sample size sufficient to guarantee
that, with probability at least \(98\%\),
\[
|\widehat p-p|<0.03,
\]
regardless of the unknown value of \(p\).

\end{problem}

\begin{solution}

From the general Bernoulli estimation bound,
\[
n
\geq
\frac1{4\varepsilon^2\alpha}
\]
is sufficient to guarantee
\[
P(|\widehat p-p|<\varepsilon)
\geq
1-\alpha.
\]

Here,
\[
\varepsilon=0.03
\]
and
\[
1-\alpha=0.98,
\]
so
\[
\alpha=0.02.
\]

Therefore,
\[
\begin{aligned}
n
&\geq
\frac{
1
}{
4(0.03)^2(0.02)
}\\
&=
13888.888\ldots.
\end{aligned}
\]

Since the sample size must be an integer,
\[
n\geq13889.
\]

Thus, at least
\[
13{,}889
\]
nails should be inspected.

\end{solution}


%%%%%%%%%%%%%%%%%%%%%%%%%%%%%%%%%%%%%%%%%%%%%%%%%%%%%%%%%%%%%%
%%%%%%%%%%%%%%%%%%%%%%  PROBLEM 56  %%%%%%%%%%%%%%%%%%%%%%%%%%
%%%%%%%%%%%%%%%%%%%%%%%%%%%%%%%%%%%%%%%%%%%%%%%%%%%%%%%%%%%%%%

\newpage
\begin{problem}

A coin has an unknown probability \(p\) of landing heads.

The coin is flipped independently \(3000\) times, and
\[
\widehat p
\]
denotes the fraction of flips that result in heads.

Use Chebyshev's inequality to show that
\[
P(|\widehat p-p|<0.03)
\geq0.90.
\]

\end{problem}

\begin{solution}

Let
\[
X_i
=
\begin{cases}
1, & \text{if the \(i\)-th flip is heads},\\
0, & \text{otherwise}.
\end{cases}
\]

Then
\[
X_i\sim\operatorname{Bernoulli}(p)
\]
and
\[
\widehat p
=
\frac1{3000}
\sum_{i=1}^{3000}X_i.
\]

Therefore,
\[
E(\widehat p)=p
\]
and
\[
\operatorname{Var}(\widehat p)
=
\frac{p(1-p)}{3000}.
\]

By Chebyshev's inequality,
\[
P(|\widehat p-p|\geq0.03)
\leq
\frac{
p(1-p)
}{
(0.03)^2(3000)
}.
\]

Since
\[
p(1-p)\leq\frac14,
\]
we obtain
\[
P(|\widehat p-p|\geq0.03)
\leq
\frac{
1
}{
4(0.03)^2(3000)
}.
\]

Thus,
\[
\begin{aligned}
P(|\widehat p-p|<0.03)
&\geq
1-
\frac{
1
}{
4(0.03)^2(3000)
}\\
&=
1-\frac1{10.8}\\
&\approx
0.9074.
\end{aligned}
\]

Therefore,
\[
P(|\widehat p-p|<0.03)
\geq
0.9074
>
0.90.
\]

Hence,
\[
P(|\widehat p-p|<0.03)
\geq0.90.
\]

\end{solution}




















%%%%%%%%%%%%%%%%%%%%%%%%%%%%%%%%%%%%%%%%%%%%%%%%%%%%%%%%%%%%%%
%%%%%%%%%%%%%%%%%%%%%%  PROBLEM 57  %%%%%%%%%%%%%%%%%%%%%%%%%%
%%%%%%%%%%%%%%%%%%%%%%%%%%%%%%%%%%%%%%%%%%%%%%%%%%%%%%%%%%%%%%

\newpage
\begin{problem}

Let \(X\) be a random variable and let \(k>0\) be a constant. Prove that
\[
P(X>t)
\leq
\frac{E(e^{kX})}{e^{kt}},
\]
whenever \(E(e^{kX})<\infty\).

\end{problem}

\begin{solution}

Since the exponential function is increasing and \(k>0\),
\[
X>t
\]
implies
\[
e^{kX}>e^{kt}.
\]

Therefore,
\[
P(X>t)
=
P(e^{kX}>e^{kt}).
\]

The random variable
\[
e^{kX}
\]
is nonnegative. Hence, by Markov's inequality,
\[
P(e^{kX}\geq e^{kt})
\leq
\frac{E(e^{kX})}{e^{kt}}.
\]

Since
\[
\{e^{kX}>e^{kt}\}
\subseteq
\{e^{kX}\geq e^{kt}\},
\]
we obtain
\[
P(X>t)
\leq
\frac{E(e^{kX})}{e^{kt}}.
\]

\end{solution}


%%%%%%%%%%%%%%%%%%%%%%%%%%%%%%%%%%%%%%%%%%%%%%%%%%%%%%%%%%%%%%
%%%%%%%%%%%%%%%%%%%%%%  PROBLEM 58  %%%%%%%%%%%%%%%%%%%%%%%%%%
%%%%%%%%%%%%%%%%%%%%%%%%%%%%%%%%%%%%%%%%%%%%%%%%%%%%%%%%%%%%%%

\newpage
\begin{problem}

Suppose that random variables \(X\) and \(Y\) satisfy
\[
E[(X-Y)^2]=0.
\]

Prove that
\[
P(X=Y)=1.
\]

\end{problem}

\begin{solution}

The random variable
\[
(X-Y)^2
\]
is nonnegative.

For every \(\varepsilon>0\), Markov's inequality gives
\[
P\left((X-Y)^2\geq\varepsilon\right)
\leq
\frac{E[(X-Y)^2]}{\varepsilon}.
\]

Since
\[
E[(X-Y)^2]=0,
\]
we obtain
\[
P\left((X-Y)^2\geq\varepsilon\right)=0
\]
for every \(\varepsilon>0\).

Now
\[
\{X\neq Y\}
=
\{(X-Y)^2>0\}.
\]

Also,
\[
\{(X-Y)^2>0\}
=
\bigcup_{n=1}^{\infty}
\left\{
(X-Y)^2\geq\frac1n
\right\}.
\]

Therefore,
\[
\begin{aligned}
P(X\neq Y)
&\leq
\sum_{n=1}^{\infty}
P\left(
(X-Y)^2\geq\frac1n
\right)\\
&=
0.
\end{aligned}
\]

Hence,
\[
P(X\neq Y)=0,
\]
and therefore
\[
P(X=Y)=1.
\]

\end{solution}


%%%%%%%%%%%%%%%%%%%%%%%%%%%%%%%%%%%%%%%%%%%%%%%%%%%%%%%%%%%%%%
%%%%%%%%%%%%%%%%%%%%%%  PROBLEM 59  %%%%%%%%%%%%%%%%%%%%%%%%%%
%%%%%%%%%%%%%%%%%%%%%%%%%%%%%%%%%%%%%%%%%%%%%%%%%%%%%%%%%%%%%%

\newpage
\begin{problem}

Let
\[
X_1,X_2,\ldots
\]
be independent and identically distributed nonnegative random variables with
density
\[
f(x)
=
\begin{cases}
4x(1-x), & 0\leq x\leq1,\\
0, & \text{otherwise}.
\end{cases}
\]

Find
\[
\lim_{n\to\infty}
\frac{X_1+\cdots+X_n}{n}.
\]

\end{problem}

\begin{solution}

As written, the proposed function is not a probability density function,
because
\[
\begin{aligned}
\int_0^1 4x(1-x)\,dx
&=
4\int_0^1(x-x^2)\,dx\\
&=
4\left(
\frac12-\frac13
\right)\\
&=
\frac23
\neq1.
\end{aligned}
\]

Therefore, the problem as stated does not define a valid distribution.

If the intended normalized density is
\[
f(x)=6x(1-x),
\qquad
0\leq x\leq1,
\]
then
\[
\begin{aligned}
E(X_i)
&=
\int_0^1
x\,6x(1-x)\,dx\\
&=
6\int_0^1(x^2-x^3)\,dx\\
&=
6\left(
\frac13-\frac14
\right)\\
&=
\frac12.
\end{aligned}
\]

By the Strong Law of Large Numbers,
\[
\frac{X_1+\cdots+X_n}{n}
\longrightarrow
E(X_1)
\]
with probability \(1\).

Thus, under the normalized density,
\[
\frac{X_1+\cdots+X_n}{n}
\longrightarrow
\frac12
\qquad
\text{almost surely}.
\]

\end{solution}


%%%%%%%%%%%%%%%%%%%%%%%%%%%%%%%%%%%%%%%%%%%%%%%%%%%%%%%%%%%%%%
%%%%%%%%%%%%%%%%%%%%%%  PROBLEM 60  %%%%%%%%%%%%%%%%%%%%%%%%%%
%%%%%%%%%%%%%%%%%%%%%%%%%%%%%%%%%%%%%%%%%%%%%%%%%%%%%%%%%%%%%%

\newpage
\begin{problem}

The time \(X\), measured in hours, required for a student to finish an
aptitude test has probability density function
\[
f(x)
=
\begin{cases}
6(x-1)(2-x),
&
1<x<2,\\
0,
&
\text{otherwise}.
\end{cases}
\]

A random sample of \(15\) students is selected.

Using the Central Limit Theorem, approximate the probability that the average
time required to complete the test is less than \(1\) hour and \(25\)
minutes.

\end{problem}

\begin{solution}

Let
\[
X_1,\ldots,X_{15}
\]
be the completion times and let
\[
\overline X
=
\frac1{15}
\sum_{i=1}^{15}X_i.
\]

First calculate the population mean:
\[
\begin{aligned}
E(X)
&=
\int_1^2
x\,6(x-1)(2-x)\,dx\\
&=
\frac32.
\end{aligned}
\]

Also,
\[
\begin{aligned}
E(X^2)
&=
\int_1^2
x^2\,6(x-1)(2-x)\,dx\\
&=
\frac{23}{10}.
\end{aligned}
\]

Hence,
\[
\begin{aligned}
\operatorname{Var}(X)
&=
E(X^2)-[E(X)]^2\\
&=
\frac{23}{10}
-
\frac94\\
&=
\frac1{20}.
\end{aligned}
\]

Thus,
\[
\mu=\frac32,
\qquad
\sigma^2=\frac1{20}.
\]

By the Central Limit Theorem,
\[
\overline X
\approx
N\left(
\frac32,
\frac{1}{20(15)}
\right)
=
N\left(
\frac32,
\frac1{300}
\right).
\]

Now
\[
1\text{ hour }25\text{ minutes}
=
1+\frac{25}{60}
=
\frac{17}{12}
\]
hours.

Therefore,
\[
\begin{aligned}
P\left(
\overline X<\frac{17}{12}
\right)
&\approx
P\left(
Z<
\frac{
\frac{17}{12}-\frac32
}{
1/\sqrt{300}
}
\right)\\
&=
P(Z<-1.443).
\end{aligned}
\]

Hence,
\[
P\left(
\overline X<\frac{17}{12}
\right)
\approx
\Phi(-1.443)
\approx
0.075.
\]

Thus, the approximate probability is
\[
0.075.
\]

\end{solution}


%%%%%%%%%%%%%%%%%%%%%%%%%%%%%%%%%%%%%%%%%%%%%%%%%%%%%%%%%%%%%%
%%%%%%%%%%%%%%%%%%%%%%  PROBLEM 61  %%%%%%%%%%%%%%%%%%%%%%%%%%
%%%%%%%%%%%%%%%%%%%%%%%%%%%%%%%%%%%%%%%%%%%%%%%%%%%%%%%%%%%%%%

\newpage
\begin{problem}

Twenty numbers are selected independently and uniformly from the interval
\((0,1)\).

Using the Central Limit Theorem, approximate the probability that their sum
is at least \(8\).

\end{problem}

\begin{solution}

Let
\[
X_1,\ldots,X_{20}
\overset{\text{i.i.d.}}{\sim}
\operatorname{Unif}(0,1)
\]
and define
\[
S_{20}
=
\sum_{i=1}^{20}X_i.
\]

For a \(\operatorname{Unif}(0,1)\) random variable,
\[
E(X_i)=\frac12
\]
and
\[
\operatorname{Var}(X_i)=\frac1{12}.
\]

Therefore,
\[
E(S_{20})
=
20\left(\frac12\right)
=
10,
\]
and
\[
\operatorname{Var}(S_{20})
=
20\left(\frac1{12}\right)
=
\frac53.
\]

By the Central Limit Theorem,
\[
S_{20}
\approx
N\left(
10,\frac53
\right).
\]

Hence,
\[
\begin{aligned}
P(S_{20}\geq8)
&\approx
P\left(
Z\geq
\frac{8-10}{\sqrt{5/3}}
\right)\\
&=
P(Z\geq-1.549).
\end{aligned}
\]

Therefore,
\[
\begin{aligned}
P(S_{20}\geq8)
&\approx
1-\Phi(-1.549)\\
&=
\Phi(1.549)\\
&\approx
0.939.
\end{aligned}
\]

Thus, the approximate probability is
\[
0.939.
\]

\end{solution}


%%%%%%%%%%%%%%%%%%%%%%%%%%%%%%%%%%%%%%%%%%%%%%%%%%%%%%%%%%%%%%
%%%%%%%%%%%%%%%%%%%%%%  PROBLEM 62  %%%%%%%%%%%%%%%%%%%%%%%%%%
%%%%%%%%%%%%%%%%%%%%%%%%%%%%%%%%%%%%%%%%%%%%%%%%%%%%%%%%%%%%%%

\newpage
\begin{problem}

A biologist wants to estimate the mean lifetime \(\ell\) of a certain type of
insect.

Suppose the insect lifetimes are independent random variables with
\[
E(X_i)=\ell
\]
and
\[
\operatorname{Var}(X_i)=1.5
\]
days squared.

The biologist estimates \(\ell\) using the sample mean
\[
\overline X.
\]

Using the Central Limit Theorem, determine approximately how large the sample
size \(n\) should be so that
\[
P(|\overline X-\ell|<0.2)
\approx
0.98.
\]

\end{problem}

\begin{solution}

For a sample of size \(n\), $E(\overline X)=\ell$ and $\operatorname{Var}(\overline X)
=
\frac{1.5}{n}.$ By the Central Limit Theorem,
\[
\frac{
\overline X-\ell
}{
\sqrt{1.5/n}
}
\approx
N(0,1).
\]

Thus,
\[
\begin{aligned}
P(|\overline X-\ell|<0.2)
&\approx
P\left(
\left|
Z
\right|
<
\frac{0.2\sqrt n}{\sqrt{1.5}}
\right).
\end{aligned}
\]

We want this probability to be approximately \(0.98\). Therefore,
\[
P(-z<Z<z)=0.98.
\]

This requires
\[
\Phi(z)=0.99.
\]

Using
\[
z_{0.99}\approx2.33,
\]
we set
\[
\frac{0.2\sqrt n}{\sqrt{1.5}}
\approx
2.33.
\]

Hence,
\[
\sqrt n
\approx
\frac{
2.33\sqrt{1.5}
}{0.2}.
\]

Squaring,
\[
n
\approx
\left(
\frac{
2.33\sqrt{1.5}
}{0.2}
\right)^2
\approx
203.58.
\]

Therefore, the sample size should be rounded up to
\[
n=204.
\]

\end{solution}


%%%%%%%%%%%%%%%%%%%%%%%%%%%%%%%%%%%%%%%%%%%%%%%%%%%%%%%%%%%%%%
%%%%%%%%%%%%%%%%%%%%%%  PROBLEM 63  %%%%%%%%%%%%%%%%%%%%%%%%%%
%%%%%%%%%%%%%%%%%%%%%%%%%%%%%%%%%%%%%%%%%%%%%%%%%%%%%%%%%%%%%%

\newpage
\begin{problem}

A random sample of size \(24\) is taken from a distribution with probability
density function
\[
f(x)
=
\begin{cases}
\dfrac19\left(x+\dfrac52\right),
&
1<x<3,\\[2mm]
0,
&
\text{otherwise}.
\end{cases}
\]

Let \(\overline X\) be the sample mean.

Using the Central Limit Theorem, approximate
\[
P(2<\overline X<2.15).
\]

\end{problem}

\begin{solution}

First calculate the population mean:
\[
\begin{aligned}
\mu
&=
E(X)\\
&=
\int_1^3
x
\frac19
\left(x+\frac52\right)
\,dx\\
&=
\frac{56}{27}.
\end{aligned}
\]

Next,
\[
\begin{aligned}
E(X^2)
&=
\int_1^3
x^2
\frac19
\left(x+\frac52\right)
\,dx\\
&=
\frac{125}{27}.
\end{aligned}
\]

Therefore,
\[
\begin{aligned}
\sigma^2
&=
E(X^2)-[E(X)]^2\\
&=
\frac{125}{27}
-
\left(\frac{56}{27}\right)^2\\
&=
\frac{239}{729}.
\end{aligned}
\]

For \(n=24\),
\[
E(\overline X)
=
\frac{56}{27},
\]
and
\[
\operatorname{Var}(\overline X)
=
\frac{239}{729(24)}.
\]

Thus,
\[
\operatorname{SD}(\overline X)
=
\sqrt{
\frac{239}{729(24)}
}.
\]

By the Central Limit Theorem,
\[
\overline X
\approx
N\left(
\frac{56}{27},
\frac{239}{729(24)}
\right).
\]

Therefore,
\[
\begin{aligned}
P(2<\overline X<2.15)
&\approx
P\left(
\frac{
2-\frac{56}{27}
}{
\sqrt{239/[729(24)]}
}
<
Z
<
\frac{
2.15-\frac{56}{27}
}{
\sqrt{239/[729(24)]}
}
\right)\\
&=
P(-0.634<Z<0.650).
\end{aligned}
\]

Hence,
\[
\begin{aligned}
P(2<\overline X<2.15)
&\approx
\Phi(0.650)-\Phi(-0.634)\\
&\approx
0.479.
\end{aligned}
\]

Thus, the approximate probability is
\[
0.479.
\]

\end{solution}


%%%%%%%%%%%%%%%%%%%%%%%%%%%%%%%%%%%%%%%%%%%%%%%%%%%%%%%%%%%%%%
%%%%%%%%%%%%%%%%%%%%%%  PROBLEM 64  %%%%%%%%%%%%%%%%%%%%%%%%%%
%%%%%%%%%%%%%%%%%%%%%%%%%%%%%%%%%%%%%%%%%%%%%%%%%%%%%%%%%%%%%%

\newpage
\begin{problem}

A random sample of size \(n\geq1\) is taken from the distribution with
probability density function
\[
f(x)
=
\frac12e^{-|x|},
\qquad
-\infty<x<\infty.
\]

Let
\[
\overline X
=
\frac1n
\sum_{i=1}^{n}X_i.
\]

Find
\[
P(\overline X>0).
\]

\end{problem}

\begin{solution}

The density satisfies
\[
f(x)=f(-x),
\]
so the distribution of each \(X_i\) is symmetric about \(0\).

Since the \(X_i\)'s are independent, their sum
\[
S_n=X_1+\cdots+X_n
\]
is also symmetric about \(0\).

Because the distribution is continuous,
\[
P(S_n=0)=0.
\]

By symmetry,
\[
P(S_n>0)
=
P(S_n<0).
\]

These two probabilities sum to \(1\), so
\[
P(S_n>0)=\frac12.
\]

Since
\[
\overline X>0
\]
if and only if
\[
S_n>0,
\]
we obtain
\[
P(\overline X>0)=\frac12.
\]

\end{solution}


%%%%%%%%%%%%%%%%%%%%%%%%%%%%%%%%%%%%%%%%%%%%%%%%%%%%%%%%%%%%%%
%%%%%%%%%%%%%%%%%%%%%%  PROBLEM 65  %%%%%%%%%%%%%%%%%%%%%%%%%%
%%%%%%%%%%%%%%%%%%%%%%%%%%%%%%%%%%%%%%%%%%%%%%%%%%%%%%%%%%%%%%

\newpage
\begin{problem}

For the scores on an achievement test given to a certain population of
students, the population mean is
\[
\mu=500
\]
and the population standard deviation is
\[
\sigma=100.
\]

Let \(\overline X\) be the mean of the scores of a random sample of
\(35\) students.

Using the Central Limit Theorem, estimate
\[
P(460<\overline X<540).
\]

\end{problem}

\begin{solution}

For a sample of size
\[
n=35,
\]
the sample mean satisfies
\[
E(\overline X)=500
\]
and
\[
\operatorname{SD}(\overline X)
=
\frac{100}{\sqrt{35}}.
\]

By the Central Limit Theorem,
\[
\overline X
\approx
N\left(
500,
\frac{100^2}{35}
\right).
\]

Therefore,
\[
\begin{aligned}
P(460<\overline X<540)
&\approx
P\left(
\frac{460-500}{100/\sqrt{35}}
<
Z
<
\frac{540-500}{100/\sqrt{35}}
\right)\\
&=
P(-2.366<Z<2.366).
\end{aligned}
\]

Hence,
\[
\begin{aligned}
P(460<\overline X<540)
&\approx
\Phi(2.366)-\Phi(-2.366)\\
&=
2\Phi(2.366)-1\\
&\approx
0.982.
\end{aligned}
\]

Thus, the approximate probability is
\[
0.982.
\]

\end{solution}












%%%%%%%%%%%%%%%%%%%%%%%%%%%%%%%%%%%%%%%%%%%%%%%%%%%%%%%%%%%%%%
%%%%%%%%%%%%%  UNIVARIATE TRANSFORMATION -- 1  %%%%%%%%%%%%%%%
%%%%%%%%%%%%%%%%%%%%%%%%%%%%%%%%%%%%%%%%%%%%%%%%%%%%%%%%%%%%%%

\newpage
\begin{problem}

Let
\[
X\sim\operatorname{Unif}(0,1),
\]
and define
\[
Y=-\ln X.
\]

\begin{enumerate}[label=(\alph*)]
    \item Find the cumulative distribution function \(F_Y(y)\).
    \item Find the probability density function \(f_Y(y)\).
    \item Identify the distribution of \(Y\).
\end{enumerate}

\end{problem}

\begin{solution}

Since
\[
0<X<1,
\]
we have
\[
-\ln X>0.
\]
Therefore,
\[
Y>0.
\]

\begin{enumerate}[label=(\alph*)]

\item For \(y\leq0\),
\[
F_Y(y)=0.
\]

Now suppose \(y>0\). Then
\[
\begin{aligned}
F_Y(y)
&=
P(Y\leq y)\\
&=
P(-\ln X\leq y).
\end{aligned}
\]

Since
\[
-\ln X\leq y
\]
is equivalent to
\[
\ln X\geq -y,
\]
and hence
\[
X\geq e^{-y},
\]
we obtain
\[
\begin{aligned}
F_Y(y)
&=
P(X\geq e^{-y})\\
&=
1-P(X<e^{-y}).
\end{aligned}
\]

Because
\[
X\sim\operatorname{Unif}(0,1),
\]
we have
\[
P(X<e^{-y})=e^{-y}.
\]

Therefore,
\[
F_Y(y)
=
1-e^{-y},
\qquad y>0.
\]

Thus,
\[
F_Y(y)
=
\begin{cases}
0, & y\leq0,\\[1mm]
1-e^{-y}, & y>0.
\end{cases}
\]

\item Differentiating the CDF,
\[
f_Y(y)
=
F_Y'(y).
\]

Hence,
\[
f_Y(y)
=
\begin{cases}
e^{-y}, & y>0,\\
0, & \text{otherwise}.
\end{cases}
\]

\item This is the density of an exponential random variable with parameter
\(1\). Therefore,
\[
Y\sim\operatorname{Exp}(1).
\]

\end{enumerate}

\end{solution}


%%%%%%%%%%%%%%%%%%%%%%%%%%%%%%%%%%%%%%%%%%%%%%%%%%%%%%%%%%%%%%
%%%%%%%%%%%%%  UNIVARIATE TRANSFORMATION -- 2  %%%%%%%%%%%%%%%
%%%%%%%%%%%%%%%%%%%%%%%%%%%%%%%%%%%%%%%%%%%%%%%%%%%%%%%%%%%%%%

\newpage
\begin{problem}

Let
\[
X\sim N(0,1),
\]
and define
\[
Y=X^2.
\]

\begin{enumerate}[label=(\alph*)]
    \item Explain why the transformation
    \[
    g(x)=x^2
    \]
    is not one-to-one.

    \item Find the cumulative distribution function \(F_Y(y)\).

    \item Find the probability density function \(f_Y(y)\).
\end{enumerate}

\end{problem}

\begin{solution}

\begin{enumerate}[label=(\alph*)]

\item The transformation
\[
g(x)=x^2
\]
is not one-to-one because, for every \(x\neq0\),
\[
g(x)=g(-x).
\]


Thus, a given positive value of \(Y\) has two possible preimages:
\[
x=\sqrt y
\qquad\text{and}\qquad
x=-\sqrt y.
\]

\item Since
\[
Y=X^2,
\]
we have
\[
Y\geq0.
\]

Therefore,
\[
F_Y(y)=0,
\qquad y<0.
\]

Now suppose \(y\geq0\). Then
\[
\begin{aligned}
F_Y(y)
&=
P(Y\leq y)\\
&=
P(X^2\leq y).
\end{aligned}
\]

The event
\[
X^2\leq y
\]
is equivalent to
\[
-\sqrt y\leq X\leq\sqrt y.
\]

Hence,
\[
\begin{aligned}
F_Y(y)
&=
P(-\sqrt y\leq X\leq\sqrt y)\\
&=
\Phi(\sqrt y)-\Phi(-\sqrt y).
\end{aligned}
\]

Using
\[
\Phi(-a)=1-\Phi(a),
\]
we obtain
\[
F_Y(y)
=
2\Phi(\sqrt y)-1.
\]

Thus,
\[
F_Y(y)
=
\begin{cases}
0, & y<0,\\[1mm]
2\Phi(\sqrt y)-1, & y\geq0.
\end{cases}
\]

\item For \(y>0\),
\[
f_Y(y)
=
\frac{d}{dy}
\left[
2\Phi(\sqrt y)-1
\right].
\]

Using the chain rule,
\[
\begin{aligned}
f_Y(y)
&=
2\phi(\sqrt y)
\frac{1}{2\sqrt y}\\
&=
\frac{\phi(\sqrt y)}{\sqrt y}.
\end{aligned}
\]

Since
\[
\phi(x)
=
\frac1{\sqrt{2\pi}}
e^{-x^2/2},
\]
we have
\[
\phi(\sqrt y)
=
\frac1{\sqrt{2\pi}}
e^{-y/2}.
\]

Therefore,
\[
f_Y(y)
=
\frac{1}{\sqrt{2\pi y}}
e^{-y/2},
\qquad y>0.
\]

Thus,
\[
f_Y(y)
=
\begin{cases}
\dfrac{1}{\sqrt{2\pi y}}e^{-y/2},
&
y>0,\\[2mm]
0,
&
\text{otherwise}.
\end{cases}
\]

\end{enumerate}

\end{solution}


%%%%%%%%%%%%%%%%%%%%%%%%%%%%%%%%%%%%%%%%%%%%%%%%%%%%%%%%%%%%%%
%%%%%%%%%%%%%%%%%%%%  CONVOLUTION -- 1  %%%%%%%%%%%%%%%%%%%%%%
%%%%%%%%%%%%%%%%%%%%%%%%%%%%%%%%%%%%%%%%%%%%%%%%%%%%%%%%%%%%%%

\newpage
\begin{problem}

Let \(X\) and \(Y\) be independent exponential random variables with common
parameter \(\lambda>0\):
\[
X,Y
\overset{\text{i.i.d.}}{\sim}
\operatorname{Exp}(\lambda).
\]

Define
\[
T=X+Y.
\]

Using the convolution formula
\[
f_T(t)
=
\int_{-\infty}^{\infty}
f_X(x)f_Y(t-x)\,dx,
\]
find the probability density function of \(T\).

Be sure to determine the correct limits of integration and the support of
\(T\).

\end{problem}

\begin{solution}

Since
\[
X,Y\sim\operatorname{Exp}(\lambda),
\]
their common density is
\[
f_X(x)=f_Y(x)
=
\begin{cases}
\lambda e^{-\lambda x}, & x>0,\\
0, & \text{otherwise}.
\end{cases}
\]

Since \(X>0\) and \(Y>0\),
\[
T=X+Y>0.
\]

Therefore,
\[
f_T(t)=0,
\qquad t\leq0.
\]

Now suppose \(t>0\). The convolution formula gives
\[
f_T(t)
=
\int_{-\infty}^{\infty}
f_X(x)f_Y(t-x)\,dx.
\]

For the integrand to be nonzero, we need both
\[
x>0
\]
and
\[
t-x>0.
\]

Thus,
\[
0<x<t.
\]

Therefore,
\[
\begin{aligned}
f_T(t)
&=
\int_0^t
\lambda e^{-\lambda x}
\lambda e^{-\lambda(t-x)}
\,dx\\
&=
\lambda^2
e^{-\lambda t}
\int_0^t dx\\
&=
\lambda^2 t e^{-\lambda t}.
\end{aligned}
\]

Hence,
\[
f_T(t)
=
\begin{cases}
\lambda^2 t e^{-\lambda t},
&
t>0,\\
0,
&
\text{otherwise}.
\end{cases}
\]
 

\end{solution}


%%%%%%%%%%%%%%%%%%%%%%%%%%%%%%%%%%%%%%%%%%%%%%%%%%%%%%%%%%%%%%
%%%%%%%%%%%%%%%%%%%%  CONVOLUTION -- 2  %%%%%%%%%%%%%%%%%%%%%%
%%%%%%%%%%%%%%%%%%%%%%%%%%%%%%%%%%%%%%%%%%%%%%%%%%%%%%%%%%%%%%

\newpage
\begin{problem}

Let
\[
X,Y
\overset{\text{i.i.d.}}{\sim}
\operatorname{Unif}(0,1),
\]
with \(X\) and \(Y\) independent, and define
\[
T=X+Y.
\]

Using the convolution formula
\[
f_T(t)
=
\int_{-\infty}^{\infty}
f_X(x)f_Y(t-x)\,dx,
\]
find the probability density function of \(T\).

\end{problem}

\begin{solution}

Since
\[
X,Y\sim\operatorname{Unif}(0,1),
\]
we have
\[
f_X(x)=f_Y(x)
=
\begin{cases}
1, & 0<x<1,\\
0, & \text{otherwise}.
\end{cases}
\]

Also,
\[
0<X<1
\qquad\text{and}\qquad
0<Y<1,
\]
so
\[
0<T<2.
\]

Therefore,
\[
f_T(t)=0
\]
outside the interval \((0,2)\).

For \(0<t\leq1\), the convolution formula gives
\[
f_T(t)
=
\int_{-\infty}^{\infty}
f_X(x)f_Y(t-x)\,dx.
\]

For the integrand to be nonzero, we need
\[
0<x<1
\]
and
\[
0<t-x<1.
\]

When \(0<t\leq1\), these conditions reduce to
\[
0<x<t.
\]

Therefore,
\[
f_T(t)
=
\int_0^t1\,dx
=
t.
\]

Now suppose
\[
1<t<2.
\]

Again, we need
\[
0<x<1
\]
and
\[
0<t-x<1.
\]

The second inequality gives
\[
t-1<x<t.
\]

Combining this with \(0<x<1\), we obtain
\[
t-1<x<1.
\]

Hence,
\[
f_T(t)
=
\int_{t-1}^{1}1\,dx
=
2-t.
\]

Therefore,
\[
f_T(t)
=
\begin{cases}
t,
&
0<t\leq1,\\[1mm]
2-t,
&
1<t<2,\\[1mm]
0,
&
\text{otherwise}.
\end{cases}
\]

\end{solution}


%%%%%%%%%%%%%%%%%%%%%%%%%%%%%%%%%%%%%%%%%%%%%%%%%%%%%%%%%%%%%%
%%%%%%%%%%%%%  PAIRWISE VS. MUTUAL INDEPENDENCE  %%%%%%%%%%%%%
%%%%%%%%%%%%%%%%%%%%%%%%%%%%%%%%%%%%%%%%%%%%%%%%%%%%%%%%%%%%%%

\newpage
\begin{problem}

Let \(X\) and \(Y\) be independent random variables satisfying
\[
P(X=1)=P(X=-1)=\frac12,
\]
and
\[
P(Y=1)=P(Y=-1)=\frac12.
\]

Define
\[
Z=XY.
\]

\begin{enumerate}[label=(\alph*)]
    \item Find the probability mass function of \(Z\).

    \item Show that \(X\) and \(Z\) are independent.

    \item Show that \(Y\) and \(Z\) are independent.

    \item Conclude that \(X,Y,Z\) are pairwise independent.

    \item Show that \(X,Y,Z\) are not mutually independent.
\end{enumerate}

\end{problem}

\begin{solution}

The four possible values of \((X,Y)\) are
\[
(1,1),\quad
(1,-1),\quad
(-1,1),\quad
(-1,-1),
\]
each with probability \(1/4\).

Since
\[
Z=XY,
\]
we obtain the table
\[
\begin{array}{c|c|c}
X & Y & Z\\
\hline
1 & 1 & 1\\
1 & -1 & -1\\
-1 & 1 & -1\\
-1 & -1 & 1
\end{array}
\]

\begin{enumerate}[label=(\alph*)]

\item We have
\[
P(Z=1)
=
P(X=Y)
=
\frac12,
\]
and
\[
P(Z=-1)
=
P(X\neq Y)
=
\frac12.
\]

Thus,
\[
p_Z(z)
=
\begin{cases}
\dfrac12,
&
z=1,-1,\\
0,
&
\text{otherwise}.
\end{cases}
\]

\item To show that \(X\) and \(Z\) are independent, consider their joint
probabilities.

For example,
\[
P(X=1,Z=1)
=
P(X=1,Y=1)
=
\frac14.
\]

Also,
\[
P(X=1)P(Z=1)
=
\frac12\cdot\frac12
=
\frac14.
\]

Similarly,
\[
P(X=1,Z=-1)
=
P(X=1,Y=-1)
=
\frac14,
\]
and
\[
P(X=1)P(Z=-1)
=
\frac14.
\]

The same calculation holds for \(X=-1\). Therefore,
\[
P(X=x,Z=z)
=
P(X=x)P(Z=z)
\]
for every
\[
x,z\in\{-1,1\}.
\]

Hence, \(X\) and \(Z\) are independent.

\item By symmetry, the same argument shows that
\[
Y
\]
and
\[
Z
\]
are independent.

Explicitly,
\[
P(Y=y,Z=z)
=
P(Y=y)P(Z=z)
\]
for all
\[
y,z\in\{-1,1\}.
\]

\item We are given that \(X\) and \(Y\) are independent. We have also shown
that \(X\) and \(Z\) are independent and that \(Y\) and \(Z\) are
independent.

Therefore, \(X,Y,Z\) are pairwise independent.

\item However,
\[
XYZ
=
XY(XY)
=
X^2Y^2
=
1
\]
with probability \(1\).

Thus, the three random variables satisfy a deterministic relationship.

For example,
\[
P(X=1,Y=1,Z=-1)=0.
\]

But
\[
P(X=1)P(Y=1)P(Z=-1)
=
\frac12\cdot\frac12\cdot\frac12
=
\frac18.
\]

Therefore,
\[
P(X=1,Y=1,Z=-1)
\neq
P(X=1)P(Y=1)P(Z=-1).
\]

Hence, \(X,Y,Z\) are not mutually independent.

\end{enumerate}

\end{solution}


%%%%%%%%%%%%%%%%%%%%%%%%%%%%%%%%%%%%%%%%%%%%%%%%%%%%%%%%%%%%%%
%%%%%%%%%%%%%%%%%%%%%  JOINT CDF  %%%%%%%%%%%%%%%%%%%%%%%%%%%%
%%%%%%%%%%%%%%%%%%%%%%%%%%%%%%%%%%%%%%%%%%%%%%%%%%%%%%%%%%%%%%

\newpage
\begin{problem}

Let \(X\) and \(Y\) have joint probability density function
\[
f(x,y)
=
\begin{cases}
2,
&
0<y<x<1,\\
0,
&
\text{otherwise}.
\end{cases}
\]

\begin{enumerate}[label=(\alph*)]
    \item Find the joint cumulative distribution function
    \[
    F_{X,Y}(x,y)
    =
    P(X\leq x,Y\leq y).
    \]

    \item Use the joint CDF to find the marginal CDFs
    \[
    F_X(x)
    \qquad\text{and}\qquad
    F_Y(y).
    \]

    \item Use the joint CDF to calculate
    \[
    P\left(
    \frac14<X\leq\frac34,\,
    \frac14<Y\leq\frac12
    \right).
    \]
\end{enumerate}

\end{problem}

\begin{solution}

The support is the triangular region
\[
0<y<x<1.
\]

\begin{enumerate}[label=(\alph*)]

\item We calculate
\[
F_{X,Y}(x,y)
=
P(X\leq x,Y\leq y).
\]

If
\[
x\leq0
\qquad\text{or}\qquad
y\leq0,
\]
then
\[
F_{X,Y}(x,y)=0.
\]

Now consider
\[
0<x<1,
\qquad
0<y<x.
\]

For a fixed value \(v\) of \(Y\), the support requires
\[
v<u<x,
\]
where \(u\) denotes the value of \(X\).

Thus,
\[
\begin{aligned}
F_{X,Y}(x,y)
&=
\int_0^y
\int_v^x
2\,du\,dv\\
&=
2\int_0^y(x-v)\,dv\\
&=
2xy-y^2.
\end{aligned}
\]

Next consider
\[
0<x<1,
\qquad
y\geq x.
\]

In this case, the condition \(Y\leq y\) does not further restrict the support
below \(X=x\). Therefore,
\[
\begin{aligned}
F_{X,Y}(x,y)
&=
P(X\leq x)\\
&=
\int_0^x
\int_0^u
2\,dv\,du\\
&=
\int_0^x2u\,du\\
&=
x^2.
\end{aligned}
\]

Now suppose
\[
x\geq1.
\]

If
\[
0<y<1,
\]
then the condition \(X\leq x\) imposes no restriction, so
\[
\begin{aligned}
F_{X,Y}(x,y)
&=
P(Y\leq y)\\
&=
\int_0^y
\int_v^1
2\,du\,dv\\
&=
2\int_0^y(1-v)\,dv\\
&=
2y-y^2.
\end{aligned}
\]

Finally, if
\[
x\geq1
\qquad\text{and}\qquad
y\geq1,
\]
then
\[
F_{X,Y}(x,y)=1.
\]

Thus,
\[
F_{X,Y}(x,y)
=
\begin{cases}
0,
&
x\leq0\text{ or }y\leq0,\\[1mm]
2xy-y^2,
&
0<y<x<1,\\[1mm]
x^2,
&
0<x<1,\ y\geq x,\\[1mm]
2y-y^2,
&
x\geq1,\ 0<y<1,\\[1mm]
1,
&
x\geq1,\ y\geq1.
\end{cases}
\]

\item The marginal CDF of \(X\) is
\[
F_X(x)
=
\lim_{y\to\infty}
F_{X,Y}(x,y).
\]

Hence,
\[
F_X(x)
=
\begin{cases}
0,
&
x\leq0,\\[1mm]
x^2,
&
0<x<1,\\[1mm]
1,
&
x\geq1.
\end{cases}
\]

Similarly,
\[
F_Y(y)
=
\lim_{x\to\infty}
F_{X,Y}(x,y).
\]

Thus,
\[
F_Y(y)
=
\begin{cases}
0,
&
y\leq0,\\[1mm]
2y-y^2,
&
0<y<1,\\[1mm]
1,
&
y\geq1.
\end{cases}
\]

\item For a joint CDF,
\[
\begin{aligned}
&P(a<X\leq b,\ c<Y\leq d)\\
&=
F_{X,Y}(b,d)
-
F_{X,Y}(a,d)
-
F_{X,Y}(b,c)
+
F_{X,Y}(a,c).
\end{aligned}
\]

Here,
\[
a=\frac14,
\qquad
b=\frac34,
\qquad
c=\frac14,
\qquad
d=\frac12.
\]

First,
\[
F_{X,Y}\left(\frac34,\frac12\right)
=
2\left(\frac34\right)\left(\frac12\right)
-
\left(\frac12\right)^2
=
\frac12.
\]

Next, since
\[
\frac12\geq\frac14,
\]
we use the \(y\geq x\) case:
\[
F_{X,Y}\left(\frac14,\frac12\right)
=
\left(\frac14\right)^2
=
\frac1{16}.
\]

Also,
\[
F_{X,Y}\left(\frac34,\frac14\right)
=
2\left(\frac34\right)\left(\frac14\right)
-
\left(\frac14\right)^2
=
\frac5{16}.
\]

Finally,
\[
F_{X,Y}\left(\frac14,\frac14\right)
=
\left(\frac14\right)^2
=
\frac1{16}.
\]

Therefore,
\[
\begin{aligned}
&P\left(
\frac14<X\leq\frac34,\,
\frac14<Y\leq\frac12
\right)\\
&=
\frac12
-
\frac1{16}
-
\frac5{16}
+
\frac1{16}\\
&=
\frac3{16}.
\end{aligned}
\]

Thus,
\[
P\left(
\frac14<X\leq\frac34,\,
\frac14<Y\leq\frac12
\right)
=
\frac3{16}.
\]

\end{enumerate}

\end{solution}




















\newpage
\section*{References}

\begin{enumerate}

    \item Sheldon Ross,
    \textit{A First Course in Probability},
    10th ed.,
    Pearson, 2019.

    \item Arman Jahangiri,
    \textit{MATH 394: Probability I --- Lecture Notes},
    Department of Mathematics,
    University of Washington,
    Summer 2026.

    \item Arman Jahangiri,
    \textit{MATH 394: Probability I --- Homework Problems and Solutions},
    Department of Mathematics,
    University of Washington,
    Summer 2026.

\end{enumerate}



\end{document}






